\section{The graphic and its central displacement}
\label{sec:graphic-central-return}

To count periodic orbits near the whole graphic, we first show that each
of them makes one passage through a fixed system of sections. Its central
height then parameterizes both halves of the singular passage. We use
this common parameter to construct the displacement and its interval of
admissible heights.

\subsection{The organizing orbit and the projective field}

At the parameter \eqref{eq:organizing-parameter}, set
\[
 H_0(x,y)=y-\frac{x^2}{2}+\frac{C_0}{2}.
\]
A direct calculation gives
\begin{equation}
\label{eq:parabola-invariance-new}
 \dot H_0=2xH_0.
\end{equation}
Thus \eqref{eq:organizing-parabola} is invariant.  On this parabola,
\begin{equation}
\label{eq:parabola-speed-new}
 \dot x=\frac{(x+1)^2+C_0-1}{2}>0.
\end{equation}
It follows that the parabola contains one regular orbit, directed from
$x=-\infty$ to $x=+\infty$.

Near the positive $y$-direction at infinity, put
\begin{equation}
\label{eq:projective-coordinates-new}
 x=\frac vz,\qquad y=\frac1z,\qquad z>0,
\end{equation}
and multiply the transformed vector field by the positive factor $z$.  We
obtain
\begin{equation}
\label{eq:projective-family-new}
\begin{aligned}
 v'&=P(v)-(1-Av+Cv^2)z,\\
 z'&=-z\eta(v,z),
\end{aligned}
\end{equation}
where
\begin{equation}
\label{eq:projective-factors-new}
\begin{aligned}
 P(v)&=\mu _1+\mu _2v-v^3,\\
 \eta(v,z)&=\mu _3+v+v^2+Cvz.
\end{aligned}
\end{equation}
The boundary drift \(q\) below therefore depends on the three
\(\mu\)-parameters, whereas the positive factor \(B\) retains the
dependence on \((A,C)\). This separation is used later when \(A\) is
varied to change the connection without changing the boundary action.
After shrinking the $(A,C)$-neighborhood of $(1,C_0)$,
$A^2<4C$ and hence
\begin{equation}
\label{eq:K-positive-new}
 1-Av+Cv^2>0\qquad(v\in\R).
\end{equation}
At the organizing parameter, $(v,z)=(0,0)$ is the nilpotent saddle--node
where the two ends of the affine orbit meet. Together with the affine
orbit, it forms the compact directed graphic $\GammaI$.

\begin{figure}[t]
\centering
% Generated by scripts/figure_organizing_graphic.py; do not edit.
% C_0 = 10, R = 3.
\newcommand{\figOneSpiral}[1]{\draw[#1] (0.016,0.000) -- (0.017,0.015) -- (0.016,0.030) -- (0.014,0.045) -- (0.010,0.059) -- (-0.003,0.067) -- (-0.013,0.056) -- (-0.018,0.042) -- (-0.022,0.027) -- (-0.025,0.013) -- (-0.027,-0.002) -- (-0.028,-0.017) -- (-0.028,-0.032) -- (-0.027,-0.047) -- (-0.025,-0.063) -- (-0.022,-0.077) -- (-0.018,-0.091) -- (-0.010,-0.104) -- (0.004,-0.108) -- (0.015,-0.099) -- (0.023,-0.086) -- (0.029,-0.072) -- (0.033,-0.058) -- (0.037,-0.042) -- (0.040,-0.027) -- (0.042,-0.012) -- (0.044,0.003) -- (0.045,0.019) -- (0.046,0.035) -- (0.046,0.051) -- (0.046,0.068) -- (0.045,0.084) -- (0.043,0.099) -- (0.041,0.114) -- (0.038,0.128) -- (0.034,0.143) -- (0.029,0.158) -- (0.022,0.171) -- (0.012,0.182) -- (-0.003,0.186) -- (-0.016,0.179) -- (-0.027,0.168) -- (-0.034,0.155) -- (-0.041,0.142) -- (-0.047,0.126) -- (-0.052,0.112) -- (-0.057,0.095) -- (-0.060,0.081) -- (-0.063,0.065) -- (-0.066,0.049) -- (-0.069,0.032) -- (-0.071,0.015) -- (-0.073,-0.003) -- (-0.074,-0.018) -- (-0.075,-0.032) -- (-0.075,-0.047) -- (-0.076,-0.062) -- (-0.076,-0.077) -- (-0.075,-0.092) -- (-0.075,-0.106) -- (-0.074,-0.121) -- (-0.072,-0.138) -- (-0.070,-0.156) -- (-0.067,-0.172) -- (-0.064,-0.188) -- (-0.061,-0.203) -- (-0.057,-0.217) -- (-0.052,-0.233) -- (-0.046,-0.247) -- (-0.039,-0.261) -- (-0.029,-0.274) -- (-0.018,-0.284) -- (-0.004,-0.289) -- (0.010,-0.287) -- (0.024,-0.280) -- (0.035,-0.270) -- (0.044,-0.259) -- (0.053,-0.245) -- (0.060,-0.232) -- (0.068,-0.216) -- (0.073,-0.202) -- (0.079,-0.188) -- (0.084,-0.172) -- (0.089,-0.155) -- (0.094,-0.136) -- (0.098,-0.122) -- (0.101,-0.107) -- (0.105,-0.092) -- (0.108,-0.076) -- (0.111,-0.059) -- (0.113,-0.043) -- (0.116,-0.025) -- (0.118,-0.007) -- (0.120,0.011) -- (0.122,0.029) -- (0.123,0.048) -- (0.125,0.067) -- (0.126,0.086) -- (0.126,0.105) -- (0.127,0.125) -- (0.127,0.144) -- (0.127,0.164) -- (0.126,0.183) -- (0.125,0.203) -- (0.124,0.222) -- (0.123,0.241) -- (0.121,0.260) -- (0.119,0.279) -- (0.117,0.297) -- (0.114,0.315) -- (0.111,0.332) -- (0.108,0.349) -- (0.104,0.365) -- (0.100,0.381) -- (0.096,0.396) -- (0.092,0.410) -- (0.087,0.424) -- (0.080,0.441) -- (0.073,0.456) -- (0.065,0.470) -- (0.057,0.483) -- (0.046,0.496) -- (0.035,0.506) -- (0.021,0.514) -- (0.006,0.518) -- (-0.009,0.518) -- (-0.025,0.513) -- (-0.038,0.505) -- (-0.051,0.494) -- (-0.062,0.483) -- (-0.072,0.469) -- (-0.083,0.454) -- (-0.091,0.441) -- (-0.098,0.427) -- (-0.106,0.412) -- (-0.113,0.395) -- (-0.121,0.378) -- (-0.128,0.359) -- (-0.135,0.340) -- (-0.142,0.319) -- (-0.146,0.305) -- (-0.150,0.290) -- (-0.154,0.275) -- (-0.158,0.259) -- (-0.162,0.243) -- (-0.166,0.227) -- (-0.169,0.210) -- (-0.173,0.193) -- (-0.176,0.176) -- (-0.179,0.158) -- (-0.182,0.141) -- (-0.185,0.123) -- (-0.188,0.104) -- (-0.190,0.086) -- (-0.192,0.067) -- (-0.194,0.048) -- (-0.196,0.029) -- (-0.198,0.010) -- (-0.199,-0.010) -- (-0.201,-0.029) -- (-0.202,-0.048) -- (-0.203,-0.068) -- (-0.203,-0.087) -- (-0.204,-0.107) -- (-0.204,-0.126) -- (-0.204,-0.146) -- (-0.204,-0.165) -- (-0.203,-0.184) -- (-0.203,-0.203) -- (-0.202,-0.222) -- (-0.201,-0.241) -- (-0.200,-0.259) -- (-0.199,-0.278) -- (-0.197,-0.296) -- (-0.195,-0.314) -- (-0.194,-0.331) -- (-0.191,-0.349) -- (-0.189,-0.366) -- (-0.187,-0.383) -- (-0.184,-0.399) -- (-0.181,-0.415) -- (-0.178,-0.431) -- (-0.175,-0.447) -- (-0.172,-0.462) -- (-0.168,-0.476) -- (-0.165,-0.491) -- (-0.161,-0.505) -- (-0.155,-0.525) -- (-0.149,-0.544) -- (-0.142,-0.563) -- (-0.135,-0.580) -- (-0.128,-0.596) -- (-0.121,-0.612) -- (-0.113,-0.626) -- (-0.105,-0.639) -- (-0.097,-0.652) -- (-0.085,-0.666) -- (-0.074,-0.679) -- (-0.062,-0.690) -- (-0.049,-0.699) -- (-0.037,-0.706) -- (-0.020,-0.711) -- (-0.004,-0.714) -- (0.014,-0.713) -- (0.031,-0.708) -- (0.045,-0.702) -- (0.060,-0.694) -- (0.074,-0.683) -- (0.086,-0.673) -- (0.097,-0.662) -- (0.108,-0.649) -- (0.119,-0.634) -- (0.131,-0.619) -- (0.142,-0.601) -- (0.150,-0.588) -- (0.157,-0.575) -- (0.165,-0.561) -- (0.172,-0.546) -- (0.180,-0.531) -- (0.188,-0.514) -- (0.195,-0.497) -- (0.203,-0.479) -- (0.210,-0.460) -- (0.217,-0.440) -- (0.225,-0.420) -- (0.232,-0.399) -- (0.239,-0.376) -- (0.246,-0.353) -- (0.253,-0.329) -- (0.259,-0.305) -- (0.266,-0.279) -- (0.273,-0.253) -- (0.279,-0.226) -- (0.285,-0.198) -- (0.288,-0.184) -- (0.291,-0.169) -- (0.294,-0.154) -- (0.296,-0.140) -- (0.299,-0.124) -- (0.302,-0.109) -- (0.304,-0.094) -- (0.307,-0.078) -- (0.310,-0.063) -- (0.312,-0.047) -- (0.314,-0.031) -- (0.317,-0.014) -- (0.319,0.002) -- (0.321,0.018) -- (0.323,0.035) -- (0.325,0.052) -- (0.327,0.069) -- (0.329,0.086) -- (0.330,0.103) -- (0.332,0.120) -- (0.333,0.138) -- (0.335,0.155) -- (0.336,0.173) -- (0.338,0.190) -- (0.339,0.208) -- (0.340,0.226) -- (0.341,0.243) -- (0.342,0.261) -- (0.342,0.279) -- (0.343,0.297) -- (0.344,0.315) -- (0.344,0.333) -- (0.344,0.351) -- (0.345,0.369) -- (0.345,0.387) -- (0.345,0.405) -- (0.345,0.422) -- (0.345,0.440) -- (0.344,0.458) -- (0.344,0.476) -- (0.344,0.494) -- (0.343,0.511) -- (0.342,0.529) -- (0.342,0.546) -- (0.341,0.563) -- (0.340,0.581) -- (0.339,0.598) -- (0.337,0.615) -- (0.336,0.632) -- (0.335,0.649) -- (0.333,0.665) -- (0.331,0.682) -- (0.330,0.698) -- (0.328,0.714) -- (0.326,0.730) -- (0.324,0.746) -- (0.322,0.762) -- (0.320,0.777) -- (0.318,0.793) -- (0.315,0.808) -- (0.313,0.823) -- (0.310,0.838) -- (0.308,0.852) -- (0.305,0.867) -- (0.302,0.881) -- (0.296,0.909) -- (0.290,0.935) -- (0.284,0.962) -- (0.277,0.987) -- (0.270,1.011) -- (0.263,1.034) -- (0.256,1.057) -- (0.248,1.079) -- (0.240,1.099) -- (0.232,1.119) -- (0.224,1.138) -- (0.216,1.156) -- (0.207,1.173) -- (0.198,1.190) -- (0.189,1.205) -- (0.180,1.220) -- (0.171,1.233) -- (0.162,1.246) -- (0.152,1.258) -- (0.142,1.269) -- (0.128,1.285) -- (0.113,1.298) -- (0.097,1.310) -- (0.082,1.319) -- (0.066,1.328) -- (0.050,1.334) -- (0.034,1.338) -- (0.018,1.341) -- (0.001,1.342) -- (-0.015,1.341) -- (-0.032,1.339) -- (-0.050,1.335) -- (-0.067,1.328) -- (-0.085,1.320) -- (-0.102,1.310) -- (-0.120,1.298) -- (-0.132,1.289) -- (-0.144,1.279) -- (-0.157,1.268) -- (-0.169,1.256) -- (-0.181,1.243) -- (-0.194,1.229) -- (-0.206,1.214) -- (-0.219,1.197) -- (-0.231,1.180) -- (-0.244,1.162) -- (-0.256,1.142) -- (-0.269,1.122) -- (-0.281,1.100) -- (-0.294,1.077) -- (-0.306,1.053) -- (-0.318,1.027) -- (-0.325,1.014) -- (-0.331,1.000) -- (-0.337,0.986) -- (-0.343,0.972) -- (-0.349,0.958) -- (-0.355,0.943) -- (-0.361,0.928) -- (-0.367,0.912) -- (-0.373,0.896) -- (-0.379,0.880) -- (-0.385,0.864) -- (-0.390,0.847) -- (-0.396,0.830) -- (-0.402,0.812) -- (-0.407,0.794) -- (-0.413,0.776) -- (-0.418,0.758) -- (-0.423,0.739) -- (-0.429,0.720) -- (-0.434,0.701) -- (-0.439,0.681) -- (-0.444,0.661) -- (-0.449,0.641) -- (-0.454,0.620) -- (-0.458,0.599) -- (-0.463,0.578) -- (-0.467,0.557) -- (-0.472,0.535) -- (-0.476,0.513) -- (-0.480,0.491) -- (-0.484,0.469) -- (-0.488,0.447) -- (-0.491,0.424) -- (-0.495,0.401) -- (-0.498,0.378) -- (-0.502,0.355) -- (-0.505,0.331) -- (-0.508,0.308) -- (-0.511,0.284) -- (-0.513,0.260) -- (-0.516,0.236) -- (-0.518,0.212) -- (-0.520,0.188) -- (-0.522,0.164) -- (-0.524,0.140) -- (-0.526,0.116) -- (-0.527,0.091) -- (-0.528,0.067) -- (-0.530,0.043) -- (-0.531,0.019) -- (-0.531,-0.006) -- (-0.532,-0.030) -- (-0.533,-0.054) -- (-0.533,-0.078) -- (-0.533,-0.101) -- (-0.533,-0.125) -- (-0.533,-0.149) -- (-0.533,-0.172) -- (-0.532,-0.196) -- (-0.532,-0.219) -- (-0.531,-0.242) -- (-0.530,-0.265) -- (-0.529,-0.287) -- (-0.528,-0.310) -- (-0.526,-0.332) -- (-0.525,-0.354) -- (-0.523,-0.376) -- (-0.521,-0.398) -- (-0.519,-0.419) -- (-0.517,-0.440) -- (-0.515,-0.461) -- (-0.513,-0.482) -- (-0.510,-0.502) -- (-0.508,-0.522) -- (-0.505,-0.542) -- (-0.503,-0.562) -- (-0.500,-0.581) -- (-0.497,-0.600) -- (-0.494,-0.619) -- (-0.491,-0.638) -- (-0.487,-0.656) -- (-0.484,-0.674) -- (-0.481,-0.692) -- (-0.477,-0.709) -- (-0.474,-0.726) -- (-0.470,-0.743) -- (-0.466,-0.760) -- (-0.463,-0.776) -- (-0.459,-0.792) -- (-0.455,-0.808) -- (-0.451,-0.823) -- (-0.447,-0.838) -- (-0.443,-0.853) -- (-0.439,-0.868) -- (-0.435,-0.882) -- (-0.430,-0.896) -- (-0.426,-0.910) -- (-0.417,-0.937) -- (-0.409,-0.963) -- (-0.400,-0.988) -- (-0.390,-1.012) -- (-0.381,-1.035) -- (-0.372,-1.057) -- (-0.362,-1.079) -- (-0.353,-1.099) -- (-0.343,-1.118) -- (-0.334,-1.137) -- (-0.324,-1.155) -- (-0.314,-1.172) -- (-0.304,-1.189) -- (-0.294,-1.204) -- (-0.284,-1.219) -- (-0.274,-1.233) -- (-0.264,-1.247) -- (-0.254,-1.260) -- (-0.244,-1.272) -- (-0.234,-1.284) -- (-0.224,-1.295) -- (-0.213,-1.305) -- (-0.198,-1.320) -- (-0.182,-1.333) -- (-0.167,-1.345) -- (-0.151,-1.356) -- (-0.135,-1.366) -- (-0.119,-1.374) -- (-0.103,-1.382) -- (-0.087,-1.388) -- (-0.071,-1.394) -- (-0.055,-1.398) -- (-0.038,-1.401) -- (-0.021,-1.403) -- (-0.004,-1.404) -- (0.013,-1.403) -- (0.030,-1.402) -- (0.048,-1.399) -- (0.066,-1.395) -- (0.084,-1.390) -- (0.102,-1.384) -- (0.121,-1.376) -- (0.140,-1.367) -- (0.153,-1.360) -- (0.166,-1.352) -- (0.180,-1.343) -- (0.193,-1.334) -- (0.207,-1.324) -- (0.221,-1.313) -- (0.235,-1.301) -- (0.249,-1.288) -- (0.264,-1.274) -- (0.279,-1.259) -- (0.294,-1.243) -- (0.309,-1.226) -- (0.325,-1.208) -- (0.341,-1.188) -- (0.357,-1.167) -- (0.373,-1.145) -- (0.382,-1.133) -- (0.390,-1.121) -- (0.399,-1.108) -- (0.407,-1.095) -- (0.416,-1.081) -- (0.425,-1.067) -- (0.433,-1.053) -- (0.442,-1.038) -- (0.451,-1.023) -- (0.460,-1.007) -- (0.469,-0.990) -- (0.479,-0.973) -- (0.488,-0.956) -- (0.497,-0.938) -- (0.506,-0.919) -- (0.516,-0.900) -- (0.525,-0.880) -- (0.535,-0.860) -- (0.544,-0.838) -- (0.554,-0.817) -- (0.564,-0.794) -- (0.573,-0.771) -- (0.583,-0.747) -- (0.593,-0.722) -- (0.603,-0.697) -- (0.612,-0.671) -- (0.622,-0.644) -- (0.632,-0.616) -- (0.642,-0.588) -- (0.651,-0.559) -- (0.661,-0.528) -- (0.670,-0.498) -- (0.680,-0.466) -- (0.689,-0.434) -- (0.698,-0.400) -- (0.707,-0.366) -- (0.716,-0.331) -- (0.725,-0.296) -- (0.734,-0.259) -- (0.742,-0.222) -- (0.751,-0.185) -- (0.758,-0.146) -- (0.766,-0.107) -- (0.774,-0.067) -- (0.781,-0.027) -- (0.788,0.014) -- (0.794,0.055) -- (0.800,0.097) -- (0.806,0.139) -- (0.811,0.182) -- (0.816,0.225) -- (0.821,0.268) -- (0.825,0.311) -- (0.828,0.354) -- (0.832,0.397) -- (0.834,0.440) -- (0.837,0.483) -- (0.839,0.526) -- (0.840,0.569) -- (0.841,0.611) -- (0.842,0.653) -- (0.842,0.694) -- (0.841,0.735) -- (0.840,0.776) -- (0.839,0.816) -- (0.838,0.855) -- (0.836,0.893) -- (0.833,0.931) -- (0.830,0.968) -- (0.827,1.005) -- (0.824,1.040) -- (0.820,1.075) -- (0.816,1.109) -- (0.812,1.142) -- (0.807,1.174) -- (0.802,1.206) -- (0.797,1.237) -- (0.792,1.266) -- (0.787,1.295) -- (0.781,1.323) -- (0.775,1.351) -- (0.769,1.377) -- (0.763,1.403) -- (0.757,1.428) -- (0.751,1.452) -- (0.745,1.475) -- (0.738,1.498) -- (0.732,1.520) -- (0.725,1.541) -- (0.719,1.562) -- (0.712,1.582) -- (0.705,1.601) -- (0.699,1.620) -- (0.692,1.638) -- (0.685,1.656) -- (0.678,1.673) -- (0.672,1.689) -- (0.665,1.705) -- (0.658,1.720) -- (0.651,1.735) -- (0.645,1.750) -- (0.638,1.764) -- (0.631,1.777) -- (0.625,1.790) -- (0.612,1.815) -- (0.599,1.839) -- (0.586,1.861) -- (0.573,1.882) -- (0.561,1.901) -- (0.548,1.919) -- (0.536,1.937) -- (0.524,1.953) -- (0.512,1.968) -- (0.500,1.983) -- (0.489,1.996) -- (0.477,2.009) -- (0.466,2.021) -- (0.455,2.033) -- (0.444,2.044) -- (0.434,2.054) -- (0.423,2.064) -- (0.407,2.078) -- (0.392,2.091) -- (0.377,2.102) -- (0.362,2.113) -- (0.348,2.123) -- (0.334,2.133) -- (0.320,2.141) -- (0.306,2.150) -- (0.292,2.157) -- (0.279,2.164) -- (0.265,2.171) -- (0.252,2.177) -- (0.235,2.184) -- (0.218,2.191) -- (0.201,2.196) -- (0.185,2.202) -- (0.168,2.207) -- (0.152,2.211) -- (0.135,2.215) -- (0.119,2.218) -- (0.103,2.221) -- (0.087,2.223) -- (0.070,2.225) -- (0.054,2.227) -- (0.038,2.228) -- (0.021,2.228) -- (0.005,2.229) -- (-0.012,2.229) -- (-0.028,2.228) -- (-0.045,2.227) -- (-0.062,2.226) -- (-0.079,2.224) -- (-0.097,2.222) -- (-0.115,2.219) -- (-0.133,2.216) -- (-0.151,2.213) -- (-0.170,2.208) -- (-0.184,2.205) -- (-0.199,2.201) -- (-0.213,2.197) -- (-0.228,2.192) -- (-0.244,2.187) -- (-0.259,2.182) -- (-0.275,2.176) -- (-0.292,2.170) -- (-0.308,2.163) -- (-0.325,2.155) -- (-0.343,2.147) -- (-0.360,2.138) -- (-0.379,2.129) -- (-0.397,2.119) -- (-0.410,2.111) -- (-0.423,2.104) -- (-0.436,2.096) -- (-0.450,2.087) -- (-0.463,2.078) -- (-0.477,2.069) -- (-0.492,2.059) -- (-0.506,2.048) -- (-0.521,2.037) -- (-0.536,2.025) -- (-0.552,2.013) -- (-0.567,2.000) -- (-0.583,1.986) -- (-0.600,1.971) -- (-0.617,1.956) -- (-0.634,1.939) -- (-0.651,1.922) -- (-0.669,1.903) -- (-0.687,1.883) -- (-0.706,1.863) -- (-0.715,1.852) -- (-0.724,1.840) -- (-0.734,1.829) -- (-0.744,1.817) -- (-0.753,1.804) -- (-0.763,1.792) -- (-0.773,1.778) -- (-0.783,1.765) -- (-0.793,1.751) -- (-0.804,1.736) -- (-0.814,1.721) -- (-0.824,1.706) -- (-0.835,1.690) -- (-0.845,1.674) -- (-0.856,1.657) -- (-0.866,1.639) -- (-0.877,1.621) -- (-0.888,1.603) -- (-0.899,1.584) -- (-0.909,1.564) -- (-0.920,1.544) -- (-0.931,1.523) -- (-0.942,1.501) -- (-0.953,1.479) -- (-0.964,1.456) -- (-0.975,1.432) -- (-0.985,1.407) -- (-0.996,1.382) -- (-1.007,1.356) -- (-1.018,1.330) -- (-1.028,1.302) -- (-1.039,1.274) -- (-1.049,1.245) -- (-1.059,1.215) -- (-1.069,1.184) -- (-1.079,1.153) -- (-1.089,1.120) -- (-1.098,1.087) -- (-1.107,1.053) -- (-1.116,1.019) -- (-1.125,0.983) -- (-1.133,0.947) -- (-1.141,0.910) -- (-1.149,0.872) -- (-1.156,0.834) -- (-1.163,0.795) -- (-1.170,0.755) -- (-1.176,0.715) -- (-1.182,0.674) -- (-1.187,0.632) -- (-1.191,0.590) -- (-1.196,0.548) -- (-1.199,0.505) -- (-1.202,0.462) -- (-1.205,0.418) -- (-1.207,0.375) -- (-1.208,0.331) -- (-1.209,0.287) -- (-1.210,0.243) -- (-1.210,0.199) -- (-1.209,0.155) -- (-1.208,0.111) -- (-1.206,0.068) -- (-1.203,0.025) -- (-1.201,-0.018) -- (-1.197,-0.061) -- (-1.193,-0.103) -- (-1.189,-0.145) -- (-1.184,-0.187) -- (-1.179,-0.227) -- (-1.174,-0.268) -- (-1.168,-0.307) -- (-1.161,-0.346) -- (-1.155,-0.385) -- (-1.148,-0.422) -- (-1.140,-0.459) -- (-1.133,-0.496) -- (-1.125,-0.531) -- (-1.117,-0.566) -- (-1.109,-0.600) -- (-1.100,-0.634) -- (-1.091,-0.666) -- (-1.083,-0.698) -- (-1.074,-0.729) -- (-1.064,-0.760) -- (-1.055,-0.789) -- (-1.046,-0.818) -- (-1.036,-0.847) -- (-1.027,-0.874) -- (-1.017,-0.901) -- (-1.007,-0.927) -- (-0.998,-0.953) -- (-0.988,-0.977) -- (-0.978,-1.002) -- (-0.968,-1.025) -- (-0.959,-1.048) -- (-0.949,-1.070) -- (-0.939,-1.092) -- (-0.929,-1.113) -- (-0.920,-1.134) -- (-0.910,-1.154) -- (-0.900,-1.173) -- (-0.891,-1.192) -- (-0.881,-1.211) -- (-0.872,-1.229) -- (-0.862,-1.246) -- (-0.853,-1.263) -- (-0.843,-1.280) -- (-0.834,-1.296) -- (-0.825,-1.312) -- (-0.816,-1.327) -- (-0.806,-1.342) -- (-0.797,-1.356) -- (-0.788,-1.371) -- (-0.779,-1.384) -- (-0.770,-1.398) -- (-0.762,-1.411) -- (-0.753,-1.424) -- (-0.744,-1.436) -- (-0.736,-1.448) -- (-0.727,-1.460) -- (-0.710,-1.483) -- (-0.693,-1.504) -- (-0.677,-1.525) -- (-0.661,-1.544) -- (-0.645,-1.563) -- (-0.629,-1.581) -- (-0.614,-1.597) -- (-0.599,-1.614) -- (-0.584,-1.629) -- (-0.569,-1.644) -- (-0.555,-1.658) -- (-0.540,-1.671) -- (-0.526,-1.684) -- (-0.512,-1.696) -- (-0.498,-1.708) -- (-0.485,-1.719) -- (-0.472,-1.729) -- (-0.458,-1.740) -- (-0.445,-1.749) -- (-0.432,-1.759) -- (-0.419,-1.768) -- (-0.407,-1.776) -- (-0.394,-1.784) -- (-0.382,-1.792) -- (-0.363,-1.803) -- (-0.345,-1.813) -- (-0.327,-1.823) -- (-0.310,-1.832) -- (-0.292,-1.840) -- (-0.275,-1.848) -- (-0.257,-1.855) -- (-0.240,-1.862) -- (-0.223,-1.868) -- (-0.206,-1.874) -- (-0.189,-1.879) -- (-0.173,-1.884) -- (-0.156,-1.888) -- (-0.139,-1.892) -- (-0.122,-1.895) -- (-0.106,-1.898) -- (-0.089,-1.901) -- (-0.072,-1.903) -- (-0.055,-1.904) -- (-0.038,-1.906) -- (-0.021,-1.906) -- (-0.004,-1.907) -- (0.014,-1.907) -- (0.032,-1.906) -- (0.049,-1.905) -- (0.067,-1.903) -- (0.086,-1.901) -- (0.104,-1.899) -- (0.123,-1.895) -- (0.142,-1.892) -- (0.162,-1.887) -- (0.182,-1.882) -- (0.202,-1.876) -- (0.216,-1.872) -- (0.230,-1.868) -- (0.245,-1.863) -- (0.259,-1.857) -- (0.274,-1.851) -- (0.290,-1.845) -- (0.305,-1.838) -- (0.321,-1.831) -- (0.337,-1.823) -- (0.354,-1.815) -- (0.371,-1.806) -- (0.389,-1.797) -- (0.407,-1.786) -- (0.425,-1.775) -- (0.444,-1.764) -- (0.463,-1.751) -- (0.484,-1.737) -- (0.504,-1.722) -- (0.526,-1.707) -- (0.548,-1.690) -- (0.559,-1.681) -- (0.571,-1.671) -- (0.582,-1.662) -- (0.594,-1.652) -- (0.607,-1.641) -- (0.619,-1.630) -- (0.631,-1.619) -- (0.644,-1.607) -- (0.657,-1.595) -- (0.671,-1.582) -- (0.684,-1.569) -- (0.698,-1.555) -- (0.712,-1.540) -- (0.727,-1.525) -- (0.741,-1.509) -- (0.756,-1.493) -- (0.772,-1.476) -- (0.787,-1.458) -- (0.803,-1.439) -- (0.819,-1.420) -- (0.836,-1.400) -- (0.853,-1.378) -- (0.870,-1.356) -- (0.887,-1.333) -- (0.905,-1.308) -- (0.923,-1.283) -- (0.942,-1.256) -- (0.960,-1.228) -- (0.980,-1.199) -- (0.999,-1.168) -- (1.019,-1.136) -- (1.039,-1.102) -- (1.059,-1.067) -- (1.080,-1.030) -- (1.100,-0.991) -- (1.121,-0.951) -- (1.142,-0.908) -- (1.163,-0.864) -- (1.185,-0.818) -- (1.206,-0.770) -- (1.227,-0.720) -- (1.247,-0.668) -- (1.268,-0.614) -- (1.288,-0.557) -- (1.308,-0.499) -- (1.327,-0.439) -- (1.345,-0.377) -- (1.363,-0.313) -- (1.380,-0.248) -- (1.396,-0.181) -- (1.410,-0.113) -- (1.424,-0.044) -- (1.436,0.027) -- (1.447,0.097) -- (1.456,0.169) -- (1.464,0.240) -- (1.470,0.311) -- (1.474,0.382) -- (1.477,0.452) -- (1.479,0.521) -- (1.479,0.589) -- (1.478,0.656) -- (1.475,0.721) -- (1.471,0.784) -- (1.465,0.846) -- (1.459,0.906) -- (1.451,0.963) -- (1.443,1.019) -- (1.433,1.072) -- (1.423,1.124) -- (1.413,1.173) -- (1.401,1.220) -- (1.390,1.265) -- (1.378,1.308) -- (1.365,1.349) -- (1.352,1.389) -- (1.339,1.426) -- (1.326,1.461) -- (1.313,1.495) -- (1.300,1.528) -- (1.287,1.558) -- (1.274,1.588) -- (1.261,1.615) -- (1.248,1.642) -- (1.235,1.667) -- (1.222,1.691) -- (1.209,1.714) -- (1.197,1.736) -- (1.184,1.757) -- (1.172,1.777) -- (1.160,1.796) -- (1.148,1.814) -- (1.136,1.831) -- (1.125,1.848) -- (1.114,1.863) -- (1.103,1.878) -- (1.092,1.893) -- (1.081,1.907) -- (1.070,1.920) -- (1.060,1.933) -- (1.050,1.945) -- (1.040,1.957) -- (1.030,1.968) -- (1.021,1.979) -- (1.002,1.999) -- (0.984,2.018) -- (0.966,2.035) -- (0.950,2.051) -- (0.934,2.067) -- (0.918,2.081) -- (0.903,2.094) -- (0.888,2.106) -- (0.874,2.118) -- (0.861,2.128) -- (0.848,2.138) -- (0.835,2.148) -- (0.823,2.157) -- (0.811,2.165) -- (0.793,2.177) -- (0.777,2.188) -- (0.761,2.198) -- (0.746,2.208) -- (0.732,2.216) -- (0.718,2.224) -- (0.704,2.232) -- (0.691,2.239) -- (0.675,2.248) -- (0.659,2.256) -- (0.644,2.263) -- (0.630,2.270) -- (0.617,2.276) -- (0.600,2.283) -- (0.585,2.290) -- (0.570,2.296) -- (0.556,2.301) -- (0.540,2.307) -- (0.525,2.313) -- (0.510,2.318) -- (0.496,2.322) -- (0.481,2.327) -- (0.466,2.332) -- (0.452,2.336) -- (0.437,2.340) -- (0.423,2.344) -- (0.407,2.348) -- (0.393,2.352) -- (0.377,2.355) -- (0.362,2.359) -- (0.347,2.362) -- (0.332,2.365) -- (0.317,2.368) -- (0.302,2.370) -- (0.287,2.373) -- (0.273,2.375) -- (0.258,2.377) -- (0.243,2.379) -- (0.229,2.381) -- (0.214,2.383) -- (0.199,2.385) -- (0.184,2.386) -- (0.169,2.388) -- (0.154,2.389) -- (0.140,2.390) -- (0.125,2.391) -- (0.110,2.392) -- (0.095,2.393) -- (0.080,2.394) -- (0.065,2.394) -- (0.050,2.395) -- (0.035,2.395) -- (0.021,2.395) -- (0.006,2.395) -- (-0.009,2.395) -- (-0.024,2.395) -- (-0.039,2.395) -- (-0.053,2.395) -- (-0.068,2.394) -- (-0.083,2.394) -- (-0.098,2.393) -- (-0.113,2.392) -- (-0.127,2.391) -- (-0.142,2.390) -- (-0.157,2.389) -- (-0.171,2.388) -- (-0.186,2.387) -- (-0.201,2.385) -- (-0.216,2.383) -- (-0.230,2.382) -- (-0.245,2.380) -- (-0.260,2.378) -- (-0.275,2.376) -- (-0.290,2.374) -- (-0.305,2.371) -- (-0.320,2.369) -- (-0.334,2.367) -- (-0.349,2.364) -- (-0.364,2.361) -- (-0.378,2.358) -- (-0.393,2.355) -- (-0.409,2.352) -- (-0.423,2.348) -- (-0.438,2.345) -- (-0.453,2.341) -- (-0.468,2.337) -- (-0.482,2.333) -- (-0.498,2.329) -- (-0.514,2.325) -- (-0.528,2.320) -- (-0.543,2.316) -- (-0.558,2.311) -- (-0.572,2.306) -- (-0.586,2.302) -- (-0.601,2.296) -- (-0.616,2.291) -- (-0.631,2.285) -- (-0.645,2.280) -- (-0.659,2.274) -- (-0.673,2.268) -- (-0.688,2.262) -- (-0.704,2.255) -- (-0.720,2.248) -- (-0.733,2.242) -- (-0.746,2.236) -- (-0.760,2.229) -- (-0.775,2.222) -- (-0.790,2.214) -- (-0.805,2.206) -- (-0.821,2.197) -- (-0.837,2.188) -- (-0.850,2.181) -- (-0.863,2.173) -- (-0.876,2.165) -- (-0.890,2.156) -- (-0.904,2.148) -- (-0.919,2.138) -- (-0.934,2.128) -- (-0.949,2.117) -- (-0.965,2.106) -- (-0.981,2.094) -- (-0.998,2.082) -- (-1.009,2.073) -- (-1.021,2.064) -- (-1.033,2.054) -- (-1.045,2.044) -- (-1.058,2.034) -- (-1.071,2.023) -- (-1.084,2.012) -- (-1.097,2.000) -- (-1.110,1.987) -- (-1.124,1.974) -- (-1.139,1.961) -- (-1.153,1.946) -- (-1.168,1.931) -- (-1.183,1.916) -- (-1.198,1.899) -- (-1.214,1.882) -- (-1.230,1.863) -- (-1.247,1.844) -- (-1.263,1.824) -- (-1.280,1.802) -- (-1.298,1.780) -- (-1.307,1.768) -- (-1.316,1.756) -- (-1.325,1.743) -- (-1.334,1.731) -- (-1.343,1.718) -- (-1.352,1.704) -- (-1.361,1.690) -- (-1.370,1.676) -- (-1.380,1.661) -- (-1.389,1.646) -- (-1.399,1.631) -- (-1.408,1.615) -- (-1.418,1.599) -- (-1.428,1.582) -- (-1.437,1.565) -- (-1.447,1.547) -- (-1.457,1.528) -- (-1.466,1.509) -- (-1.476,1.490) -- (-1.486,1.470) -- (-1.496,1.450) -- (-1.505,1.429) -- (-1.515,1.407) -- (-1.525,1.384) -- (-1.534,1.362) -- (-1.544,1.338) -- (-1.553,1.314) -- (-1.563,1.289) -- (-1.572,1.263) -- (-1.581,1.237) -- (-1.590,1.210) -- (-1.599,1.182) -- (-1.608,1.154) -- (-1.616,1.125) -- (-1.625,1.095) -- (-1.633,1.064) -- (-1.640,1.033) -- (-1.648,1.001) -- (-1.655,0.968) -- (-1.662,0.934) -- (-1.668,0.900) -- (-1.675,0.865) -- (-1.680,0.830) -- (-1.686,0.793) -- (-1.690,0.756) -- (-1.695,0.719) -- (-1.699,0.681) -- (-1.702,0.642) -- (-1.705,0.602) -- (-1.708,0.563) -- (-1.709,0.522) -- (-1.711,0.481) -- (-1.711,0.440) -- (-1.711,0.399) -- (-1.711,0.357) -- (-1.710,0.315) -- (-1.708,0.272) -- (-1.706,0.230) -- (-1.703,0.187) -- (-1.699,0.145) -- (-1.695,0.102) -- (-1.690,0.060) -- (-1.685,0.017) -- (-1.679,-0.025) -- (-1.672,-0.067) -- (-1.665,-0.109) -- (-1.657,-0.151) -- (-1.649,-0.192) -- (-1.640,-0.232) -- (-1.631,-0.273) -- (-1.622,-0.313) -- (-1.612,-0.352) -- (-1.601,-0.391) -- (-1.590,-0.429) -- (-1.579,-0.467) -- (-1.568,-0.504) -- (-1.556,-0.540) -- (-1.544,-0.576) -- (-1.532,-0.611) -- (-1.519,-0.645) -- (-1.506,-0.679) -- (-1.494,-0.712) -- (-1.480,-0.744) -- (-1.467,-0.776) -- (-1.454,-0.806) -- (-1.441,-0.837) -- (-1.427,-0.866) -- (-1.413,-0.895) -- (-1.400,-0.923) -- (-1.386,-0.950) -- (-1.373,-0.977) -- (-1.359,-1.003) -- (-1.345,-1.028) -- (-1.332,-1.053) -- (-1.318,-1.077) -- (-1.304,-1.101) -- (-1.291,-1.124) -- (-1.278,-1.146) -- (-1.264,-1.168) -- (-1.251,-1.189) -- (-1.238,-1.209) -- (-1.224,-1.229) -- (-1.211,-1.249) -- (-1.198,-1.268) -- (-1.186,-1.287) -- (-1.173,-1.305) -- (-1.160,-1.322) -- (-1.148,-1.339) -- (-1.135,-1.356) -- (-1.123,-1.372) -- (-1.111,-1.388) -- (-1.099,-1.403) -- (-1.087,-1.418) -- (-1.075,-1.433) -- (-1.063,-1.447) -- (-1.052,-1.461) -- (-1.040,-1.475) -- (-1.029,-1.488) -- (-1.017,-1.501) -- (-1.006,-1.513) -- (-0.995,-1.525) -- (-0.984,-1.537) -- (-0.974,-1.549) -- (-0.963,-1.560) -- (-0.952,-1.571) -- (-0.942,-1.582) -- (-0.932,-1.593) -- (-0.921,-1.603) -- (-0.901,-1.623) -- (-0.881,-1.641) -- (-0.862,-1.659) -- (-0.843,-1.677) -- (-0.825,-1.693) -- (-0.807,-1.709) -- (-0.789,-1.724) -- (-0.771,-1.738) -- (-0.754,-1.752) -- (-0.737,-1.765) -- (-0.721,-1.777) -- (-0.705,-1.789) -- (-0.689,-1.801) -- (-0.673,-1.812) -- (-0.658,-1.822) -- (-0.643,-1.832) -- (-0.628,-1.842) -- (-0.613,-1.852) -- (-0.599,-1.861) -- (-0.585,-1.869) -- (-0.571,-1.878) -- (-0.557,-1.885) -- (-0.544,-1.893) -- (-0.531,-1.901) -- (-0.517,-1.908) -- (-0.504,-1.914) -- (-0.485,-1.924) -- (-0.467,-1.933) -- (-0.448,-1.942) -- (-0.430,-1.950) -- (-0.412,-1.958) -- (-0.394,-1.965) -- (-0.377,-1.972) -- (-0.360,-1.979) -- (-0.343,-1.985) -- (-0.326,-1.991) -- (-0.310,-1.996) -- (-0.293,-2.001) -- (-0.277,-2.006) -- (-0.261,-2.010) -- (-0.245,-2.015) -- (-0.229,-2.019) -- (-0.213,-2.022) -- (-0.198,-2.026) -- (-0.182,-2.029) -- (-0.166,-2.031) -- (-0.151,-2.034) -- (-0.135,-2.036) -- (-0.120,-2.038) -- (-0.104,-2.040) -- (-0.089,-2.042) -- (-0.073,-2.043) -- (-0.057,-2.044) -- (-0.042,-2.045) -- (-0.026,-2.045) -- (-0.010,-2.046) -- (0.006,-2.046) -- (0.022,-2.046) -- (0.039,-2.045) -- (0.055,-2.044) -- (0.072,-2.043) -- (0.089,-2.042) -- (0.106,-2.040) -- (0.123,-2.038) -- (0.141,-2.036) -- (0.159,-2.033) -- (0.177,-2.030) -- (0.196,-2.026) -- (0.215,-2.022) -- (0.234,-2.017) -- (0.254,-2.012) -- (0.275,-2.007) -- (0.289,-2.003) -- (0.303,-1.999) -- (0.318,-1.994) -- (0.332,-1.989) -- (0.348,-1.984) -- (0.363,-1.978) -- (0.379,-1.972) -- (0.395,-1.966) -- (0.412,-1.959) -- (0.429,-1.952) -- (0.446,-1.944) -- (0.464,-1.936) -- (0.483,-1.927) -- (0.502,-1.917) -- (0.521,-1.907) -- (0.541,-1.896) -- (0.562,-1.884) -- (0.584,-1.872) -- (0.606,-1.858) -- (0.629,-1.844) -- (0.653,-1.828) -- (0.666,-1.820) -- (0.678,-1.811) -- (0.691,-1.802) -- (0.704,-1.793) -- (0.717,-1.783) -- (0.731,-1.773) -- (0.745,-1.763) -- (0.759,-1.752) -- (0.773,-1.741) -- (0.788,-1.729) -- (0.803,-1.716) -- (0.819,-1.704) -- (0.834,-1.690) -- (0.850,-1.676) -- (0.867,-1.661) -- (0.884,-1.646) -- (0.901,-1.630) -- (0.919,-1.613) -- (0.937,-1.596) -- (0.955,-1.577) -- (0.974,-1.558) -- (0.993,-1.538) -- (1.013,-1.516) -- (1.033,-1.494) -- (1.054,-1.471) -- (1.075,-1.446) -- (1.096,-1.420) -- (1.118,-1.393) -- (1.141,-1.364) -- (1.164,-1.334) -- (1.187,-1.303) -- (1.211,-1.269) -- (1.235,-1.234) -- (1.260,-1.197) -- (1.285,-1.159) -- (1.310,-1.118) -- (1.336,-1.075) -- (1.362,-1.030) -- (1.388,-0.982) -- (1.414,-0.933) -- (1.440,-0.881) -- (1.466,-0.826) -- (1.491,-0.769) -- (1.517,-0.710) -- (1.541,-0.648) -- (1.565,-0.583) -- (1.589,-0.517) -- (1.611,-0.448) -- (1.632,-0.377) -- (1.652,-0.304) -- (1.670,-0.229) -- (1.686,-0.153) -- (1.701,-0.076) -- (1.714,0.002) -- (1.725,0.081) -- (1.734,0.160) -- (1.740,0.238) -- (1.745,0.316) -- (1.747,0.393) -- (1.748,0.469) -- (1.746,0.543) -- (1.742,0.616) -- (1.737,0.687) -- (1.730,0.755) -- (1.721,0.822) -- (1.712,0.885) -- (1.700,0.947) -- (1.688,1.006) -- (1.675,1.062) -- (1.661,1.116) -- (1.647,1.167) -- (1.632,1.216) -- (1.616,1.262) -- (1.600,1.306) -- (1.584,1.348) -- (1.568,1.388) -- (1.551,1.426) -- (1.535,1.462) -- (1.518,1.496) -- (1.502,1.528) -- (1.486,1.559) -- (1.470,1.588) -- (1.454,1.615) -- (1.438,1.642) -- (1.423,1.666) -- (1.408,1.690) -- (1.392,1.713) -- (1.378,1.734) -- (1.363,1.754) -- (1.349,1.774) -- (1.335,1.792) -- (1.321,1.810) -- (1.308,1.827) -- (1.295,1.843) -- (1.282,1.858) -- (1.269,1.873) -- (1.257,1.887) -- (1.245,1.900) -- (1.233,1.913) -- (1.221,1.925) -- (1.210,1.937) -- (1.199,1.948) -- (1.188,1.959) -- (1.177,1.969) -- (1.167,1.979) -- (1.146,1.998) -- (1.127,2.016) -- (1.108,2.032) -- (1.090,2.047) -- (1.073,2.061) -- (1.056,2.074) -- (1.040,2.086) -- (1.025,2.098) -- (1.010,2.108) -- (0.996,2.118) -- (0.982,2.128) -- (0.969,2.137) -- (0.956,2.145) -- (0.943,2.153) -- (0.926,2.164) -- (0.909,2.175) -- (0.892,2.184) -- (0.877,2.193) -- (0.862,2.201) -- (0.848,2.209) -- (0.834,2.216) -- (0.821,2.223) -- (0.804,2.231) -- (0.789,2.239) -- (0.774,2.246) -- (0.759,2.252) -- (0.746,2.258) -- (0.729,2.265) -- (0.714,2.272) -- (0.700,2.278) -- (0.686,2.283) -- (0.670,2.289) -- (0.655,2.294) -- (0.641,2.299) -- (0.626,2.305) -- (0.611,2.310) -- (0.598,2.314) -- (0.583,2.319) -- (0.569,2.323) -- (0.554,2.327) -- (0.539,2.332) -- (0.525,2.335) -- (0.510,2.339) -- (0.495,2.343) -- (0.480,2.347) -- (0.466,2.350) -- (0.452,2.353) -- (0.437,2.357) -- (0.422,2.360) -- (0.408,2.362) -- (0.393,2.365) -- (0.379,2.368) -- (0.364,2.370) -- (0.350,2.373) -- (0.335,2.375) -- (0.321,2.377) -- (0.307,2.379) -- (0.292,2.381) -- (0.278,2.383) -- (0.263,2.385) -- (0.249,2.387) -- (0.234,2.388) -- (0.220,2.390) -- (0.205,2.391) -- (0.191,2.392) -- (0.176,2.393) -- (0.162,2.394) -- (0.148,2.395) -- (0.133,2.396) -- (0.119,2.397) -- (0.104,2.398) -- (0.090,2.398) -- (0.075,2.399) -- (0.061,2.399) -- (0.047,2.400) -- (0.032,2.400) -- (0.018,2.400) -- (0.003,2.400) -- (-0.011,2.400) -- (-0.025,2.400) -- (-0.040,2.400) -- (-0.054,2.399) -- (-0.069,2.399) -- (-0.083,2.399) -- (-0.098,2.398) -- (-0.112,2.397) -- (-0.126,2.397) -- (-0.141,2.396) -- (-0.155,2.395) -- (-0.170,2.394) -- (-0.184,2.393) -- (-0.199,2.392) -- (-0.213,2.390) -- (-0.227,2.389) -- (-0.242,2.387) -- (-0.256,2.386) -- (-0.271,2.384) -- (-0.285,2.382) -- (-0.299,2.380) -- (-0.314,2.379) -- (-0.328,2.376) -- (-0.343,2.374) -- (-0.357,2.372) -- (-0.371,2.369) -- (-0.386,2.367) -- (-0.400,2.364) -- (-0.415,2.361) -- (-0.429,2.359) -- (-0.444,2.355) -- (-0.458,2.352) -- (-0.473,2.349) -- (-0.488,2.346) -- (-0.502,2.342) -- (-0.516,2.339) -- (-0.530,2.335) -- (-0.544,2.331) -- (-0.559,2.327) -- (-0.573,2.323) -- (-0.587,2.318) -- (-0.602,2.314) -- (-0.616,2.309) -- (-0.632,2.304) -- (-0.646,2.299) -- (-0.660,2.294) -- (-0.674,2.289) -- (-0.689,2.283) -- (-0.703,2.278) -- (-0.718,2.272) -- (-0.731,2.267) -- (-0.745,2.261) -- (-0.760,2.254) -- (-0.775,2.248) -- (-0.789,2.241) -- (-0.803,2.234) -- (-0.818,2.227) -- (-0.831,2.221) -- (-0.845,2.214) -- (-0.859,2.207) -- (-0.873,2.199) -- (-0.888,2.190) -- (-0.901,2.183) -- (-0.914,2.176) -- (-0.927,2.168) -- (-0.941,2.159) -- (-0.956,2.150) -- (-0.970,2.141) -- (-0.986,2.131) -- (-0.998,2.123) -- (-1.010,2.114) -- (-1.023,2.105) -- (-1.036,2.096) -- (-1.049,2.086) -- (-1.063,2.076) -- (-1.078,2.065) -- (-1.092,2.053) -- (-1.108,2.041) -- (-1.123,2.027) -- (-1.139,2.014) -- (-1.150,2.004) -- (-1.162,1.994) -- (-1.173,1.983) -- (-1.185,1.972) -- (-1.197,1.961) -- (-1.209,1.949) -- (-1.222,1.936) -- (-1.235,1.923) -- (-1.248,1.909) -- (-1.262,1.895) -- (-1.275,1.880) -- (-1.289,1.864) -- (-1.304,1.847) -- (-1.318,1.830) -- (-1.333,1.812) -- (-1.348,1.792) -- (-1.364,1.772) -- (-1.380,1.751) -- (-1.396,1.729) -- (-1.412,1.705) -- (-1.420,1.693) -- (-1.429,1.680) -- (-1.437,1.667) -- (-1.446,1.654) -- (-1.454,1.641) -- (-1.463,1.627) -- (-1.471,1.612) -- (-1.480,1.597) -- (-1.489,1.582) -- (-1.498,1.567) -- (-1.507,1.551) -- (-1.515,1.534) -- (-1.524,1.517) -- (-1.533,1.500) -- (-1.542,1.482) -- (-1.551,1.464) -- (-1.560,1.445) -- (-1.569,1.426) -- (-1.578,1.406) -- (-1.587,1.385) -- (-1.596,1.364) -- (-1.605,1.343) -- (-1.613,1.320) -- (-1.622,1.298) -- (-1.631,1.274) -- (-1.639,1.250) -- (-1.648,1.226) -- (-1.656,1.200) -- (-1.664,1.174) -- (-1.672,1.148) -- (-1.680,1.121) -- (-1.688,1.093) -- (-1.695,1.064) -- (-1.703,1.035) -- (-1.710,1.005) -- (-1.716,0.974) -- (-1.723,0.943) -- (-1.729,0.911) -- (-1.735,0.878) -- (-1.740,0.844) -- (-1.745,0.810) -- (-1.750,0.776) -- (-1.754,0.740) -- (-1.758,0.705) -- (-1.762,0.668) -- (-1.765,0.631) -- (-1.767,0.593) -- (-1.769,0.555) -- (-1.771,0.517) -- (-1.772,0.478) -- (-1.772,0.438) -- (-1.772,0.398) -- (-1.771,0.358) -- (-1.770,0.318) -- (-1.768,0.277) -- (-1.765,0.236) -- (-1.762,0.195) -- (-1.759,0.154) -- (-1.754,0.113) -- (-1.749,0.072) -- (-1.744,0.030) -- (-1.738,-0.011) -- (-1.731,-0.051) -- (-1.724,-0.092) -- (-1.717,-0.133) -- (-1.709,-0.173) -- (-1.700,-0.213) -- (-1.691,-0.252) -- (-1.681,-0.291) -- (-1.671,-0.330) -- (-1.661,-0.368) -- (-1.650,-0.406) -- (-1.639,-0.443) -- (-1.628,-0.480) -- (-1.616,-0.516) -- (-1.604,-0.551) -- (-1.591,-0.586) -- (-1.579,-0.620) -- (-1.566,-0.654) -- (-1.553,-0.687) -- (-1.540,-0.719) -- (-1.526,-0.751) -- (-1.513,-0.782) -- (-1.499,-0.812) -- (-1.485,-0.842) -- (-1.472,-0.871) -- (-1.458,-0.899) -- (-1.444,-0.927) -- (-1.430,-0.954) -- (-1.416,-0.980) -- (-1.402,-1.006) -- (-1.388,-1.031) -- (-1.374,-1.056) -- (-1.360,-1.080) -- (-1.346,-1.103) -- (-1.333,-1.126) -- (-1.319,-1.148) -- (-1.305,-1.169) -- (-1.292,-1.190) -- (-1.278,-1.211) -- (-1.265,-1.231) -- (-1.251,-1.250) -- (-1.238,-1.269) -- (-1.225,-1.288) -- (-1.212,-1.306) -- (-1.199,-1.323) -- (-1.186,-1.340) -- (-1.173,-1.357) -- (-1.161,-1.373) -- (-1.148,-1.389) -- (-1.136,-1.405) -- (-1.124,-1.420) -- (-1.112,-1.434) -- (-1.100,-1.448) -- (-1.088,-1.462) -- (-1.076,-1.476) -- (-1.064,-1.489) -- (-1.053,-1.502) -- (-1.041,-1.515) -- (-1.030,-1.527) -- (-1.019,-1.539) -- (-1.008,-1.550) -- (-0.997,-1.562) -- (-0.986,-1.573) -- (-0.975,-1.584) -- (-0.965,-1.594) -- (-0.954,-1.605) -- (-0.944,-1.615) -- (-0.924,-1.634) -- (-0.904,-1.653) -- (-0.884,-1.670) -- (-0.865,-1.687) -- (-0.846,-1.703) -- (-0.828,-1.719) -- (-0.810,-1.734) -- (-0.792,-1.748) -- (-0.775,-1.761) -- (-0.758,-1.774) -- (-0.741,-1.787) -- (-0.725,-1.798) -- (-0.709,-1.810) -- (-0.693,-1.821) -- (-0.677,-1.831) -- (-0.662,-1.841) -- (-0.647,-1.851) -- (-0.632,-1.860) -- (-0.618,-1.869) -- (-0.604,-1.878) -- (-0.590,-1.886) -- (-0.576,-1.894) -- (-0.562,-1.901) -- (-0.549,-1.909) -- (-0.535,-1.916) -- (-0.522,-1.923) -- (-0.509,-1.929) -- (-0.490,-1.939) -- (-0.472,-1.947) -- (-0.453,-1.956) -- (-0.435,-1.964) -- (-0.418,-1.971) -- (-0.400,-1.978) -- (-0.383,-1.985) -- (-0.366,-1.991) -- (-0.349,-1.997) -- (-0.332,-2.003) -- (-0.316,-2.008) -- (-0.300,-2.013) -- (-0.283,-2.018) -- (-0.267,-2.022) -- (-0.252,-2.026) -- (-0.236,-2.030) -- (-0.220,-2.033) -- (-0.205,-2.037) -- (-0.189,-2.040) -- (-0.174,-2.043) -- (-0.158,-2.045) -- (-0.143,-2.047) -- (-0.127,-2.050) -- (-0.112,-2.051) -- (-0.096,-2.053) -- (-0.081,-2.054) -- (-0.066,-2.056) -- (-0.050,-2.056) -- (-0.034,-2.057) -- (-0.019,-2.057) -- (-0.003,-2.058) -- (0.013,-2.058) -- (0.029,-2.057) -- (0.045,-2.057) -- (0.062,-2.056) -- (0.078,-2.055) -- (0.095,-2.053) -- (0.112,-2.051) -- (0.129,-2.049) -- (0.147,-2.047) -- (0.165,-2.044) -- (0.183,-2.041) -- (0.202,-2.037) -- (0.221,-2.033) -- (0.240,-2.029) -- (0.260,-2.024) -- (0.281,-2.018) -- (0.294,-2.014) -- (0.309,-2.010) -- (0.323,-2.006) -- (0.338,-2.001) -- (0.353,-1.996) -- (0.368,-1.990) -- (0.384,-1.984) -- (0.400,-1.978) -- (0.417,-1.971) -- (0.434,-1.964) -- (0.451,-1.957) -- (0.469,-1.949) -- (0.488,-1.940) -- (0.507,-1.931) -- (0.526,-1.921) -- (0.546,-1.910) -- (0.567,-1.899) -- (0.589,-1.886) -- (0.611,-1.873) -- (0.634,-1.859) -- (0.658,-1.844) -- (0.670,-1.836) -- (0.683,-1.828) -- (0.695,-1.819) -- (0.708,-1.810) -- (0.722,-1.801) -- (0.735,-1.791) -- (0.749,-1.781) -- (0.763,-1.770) -- (0.778,-1.759) -- (0.792,-1.748) -- (0.807,-1.736) -- (0.823,-1.723) -- (0.838,-1.710) -- (0.854,-1.696) -- (0.871,-1.682) -- (0.888,-1.667) -- (0.905,-1.652) -- (0.922,-1.635) -- (0.940,-1.618) -- (0.959,-1.600) -- (0.978,-1.582) -- (0.997,-1.562) -- (1.017,-1.542) -- (1.037,-1.520) -- (1.057,-1.497) -- (1.078,-1.473) -- (1.100,-1.448) -- (1.122,-1.422) -- (1.145,-1.394) -- (1.168,-1.365) -- (1.191,-1.334) -- (1.215,-1.302) -- (1.239,-1.268) -- (1.264,-1.232) -- (1.289,-1.195) -- (1.315,-1.155) -- (1.341,-1.113) -- (1.367,-1.070) -- (1.393,-1.023) -- (1.419,-0.975) -- (1.446,-0.924) -- (1.472,-0.871) -- (1.498,-0.815) -- (1.524,-0.757) -- (1.550,-0.696) -- (1.575,-0.633) -- (1.599,-0.567) -- (1.622,-0.500) -- (1.644,-0.429) -- (1.665,-0.357) -- (1.684,-0.283) -- (1.702,-0.207) -- (1.718,-0.130) -- (1.732,-0.052) -- (1.744,0.027) -- (1.754,0.106) -- (1.762,0.185) -- (1.768,0.263) -- (1.772,0.342) -- (1.773,0.419) -- (1.772,0.494) -- (1.770,0.568) -- (1.765,0.641) -- (1.759,0.711) -- (1.751,0.779) -- (1.741,0.844) -- (1.731,0.907) -- (1.719,0.968) -- (1.706,1.025) -- (1.692,1.081) -- (1.678,1.134) -- (1.662,1.184) -- (1.647,1.232) -- (1.631,1.278) -- (1.614,1.321) -- (1.598,1.362) -- (1.581,1.401) -- (1.565,1.438) -- (1.548,1.473) -- (1.531,1.506) -- (1.515,1.538) -- (1.498,1.568) -- (1.482,1.596) -- (1.466,1.623) -- (1.450,1.649) -- (1.435,1.673) -- (1.419,1.696) -- (1.404,1.719) -- (1.389,1.739) -- (1.375,1.759) -- (1.360,1.778) -- (1.346,1.796) -- (1.333,1.814) -- (1.319,1.830) -- (1.306,1.846) -- (1.293,1.861) -- (1.280,1.875) -- (1.268,1.889) -- (1.256,1.902) -- (1.244,1.915) -- (1.232,1.927) -- (1.221,1.938) -- (1.210,1.949) -- (1.199,1.960) -- (1.188,1.970) -- (1.168,1.990) -- (1.148,2.007) -- (1.129,2.024) -- (1.110,2.039) -- (1.093,2.054) -- (1.076,2.067) -- (1.059,2.079) -- (1.044,2.091) -- (1.029,2.102) -- (1.014,2.112) -- (1.000,2.122) -- (0.987,2.131) -- (0.973,2.140) -- (0.961,2.148) -- (0.949,2.156) -- (0.931,2.166) -- (0.914,2.176) -- (0.898,2.186) -- (0.883,2.194) -- (0.868,2.202) -- (0.854,2.210) -- (0.841,2.217) -- (0.828,2.223) -- (0.812,2.231) -- (0.796,2.239) -- (0.781,2.245) -- (0.767,2.252) -- (0.754,2.258) -- (0.738,2.265) -- (0.722,2.271) -- (0.708,2.277) -- (0.695,2.282) -- (0.679,2.288) -- (0.664,2.293) -- (0.651,2.298) -- (0.635,2.303) -- (0.621,2.308) -- (0.606,2.313) -- (0.591,2.318) -- (0.576,2.322) -- (0.561,2.326) -- (0.547,2.331) -- (0.533,2.334) -- (0.518,2.338) -- (0.504,2.342) -- (0.489,2.346) -- (0.474,2.349) -- (0.460,2.352) -- (0.445,2.355) -- (0.431,2.358) -- (0.416,2.361) -- (0.402,2.364) -- (0.388,2.367) -- (0.373,2.369) -- (0.359,2.372) -- (0.345,2.374) -- (0.330,2.376) -- (0.316,2.378) -- (0.302,2.380) -- (0.287,2.382) -- (0.272,2.384) -- (0.258,2.386) -- (0.244,2.387) -- (0.229,2.389) -- (0.215,2.390) -- (0.200,2.391) -- (0.186,2.393) -- (0.172,2.394) -- (0.157,2.395) -- (0.143,2.396) -- (0.128,2.397) -- (0.114,2.397) -- (0.099,2.398) -- (0.085,2.398) -- (0.071,2.399) -- (0.056,2.399) -- (0.042,2.400) -- (0.027,2.400) -- (0.013,2.400) -- (-0.001,2.400) -- (-0.016,2.400) -- (-0.019,2.400);}
\newcommand{\figOneGraphic}[1]{\draw[#1] (-0.000,2.400) -- (-0.010,2.400) -- (-0.019,2.400) -- (-0.029,2.400) -- (-0.039,2.400) -- (-0.048,2.400) -- (-0.058,2.399) -- (-0.068,2.399) -- (-0.078,2.399) -- (-0.087,2.398) -- (-0.097,2.398) -- (-0.107,2.398) -- (-0.117,2.397) -- (-0.126,2.397) -- (-0.136,2.396) -- (-0.146,2.396) -- (-0.156,2.395) -- (-0.166,2.394) -- (-0.175,2.394) -- (-0.185,2.393) -- (-0.195,2.392) -- (-0.205,2.391) -- (-0.215,2.390) -- (-0.224,2.389) -- (-0.234,2.388) -- (-0.244,2.387) -- (-0.254,2.386) -- (-0.264,2.385) -- (-0.274,2.384) -- (-0.284,2.383) -- (-0.294,2.381) -- (-0.304,2.380) -- (-0.313,2.379) -- (-0.323,2.377) -- (-0.333,2.376) -- (-0.343,2.374) -- (-0.353,2.373) -- (-0.363,2.371) -- (-0.373,2.369) -- (-0.382,2.368) -- (-0.392,2.366) -- (-0.402,2.364) -- (-0.412,2.362) -- (-0.421,2.360) -- (-0.431,2.358) -- (-0.441,2.356) -- (-0.451,2.354) -- (-0.460,2.352) -- (-0.470,2.350) -- (-0.480,2.348) -- (-0.490,2.345) -- (-0.499,2.343) -- (-0.509,2.341) -- (-0.519,2.338) -- (-0.529,2.336) -- (-0.538,2.333) -- (-0.548,2.330) -- (-0.557,2.328) -- (-0.568,2.325) -- (-0.578,2.322) -- (-0.587,2.319) -- (-0.597,2.316) -- (-0.606,2.313) -- (-0.616,2.310) -- (-0.626,2.306) -- (-0.636,2.303) -- (-0.645,2.300) -- (-0.654,2.297) -- (-0.664,2.293) -- (-0.674,2.290) -- (-0.684,2.286) -- (-0.694,2.282) -- (-0.703,2.279) -- (-0.712,2.275) -- (-0.721,2.271) -- (-0.730,2.268) -- (-0.740,2.264) -- (-0.750,2.259) -- (-0.760,2.255) -- (-0.770,2.251) -- (-0.780,2.246) -- (-0.789,2.242) -- (-0.798,2.238) -- (-0.807,2.233) -- (-0.816,2.229) -- (-0.826,2.224) -- (-0.835,2.220) -- (-0.845,2.215) -- (-0.855,2.210) -- (-0.865,2.204) -- (-0.875,2.199) -- (-0.883,2.194) -- (-0.892,2.189) -- (-0.901,2.184) -- (-0.909,2.179) -- (-0.918,2.174) -- (-0.927,2.169) -- (-0.937,2.163) -- (-0.946,2.157) -- (-0.955,2.152) -- (-0.965,2.146) -- (-0.974,2.139) -- (-0.984,2.133) -- (-0.994,2.126) -- (-1.002,2.121) -- (-1.010,2.115) -- (-1.018,2.110) -- (-1.026,2.104) -- (-1.035,2.098) -- (-1.043,2.092) -- (-1.052,2.086) -- (-1.060,2.079) -- (-1.069,2.073) -- (-1.077,2.066) -- (-1.086,2.059) -- (-1.095,2.052) -- (-1.104,2.045) -- (-1.113,2.037) -- (-1.122,2.030) -- (-1.131,2.022) -- (-1.141,2.014) -- (-1.150,2.006) -- (-1.159,1.997) -- (-1.167,1.991) -- (-1.174,1.984) -- (-1.181,1.977) -- (-1.188,1.971) -- (-1.196,1.964) -- (-1.203,1.957) -- (-1.210,1.949) -- (-1.218,1.942) -- (-1.225,1.934) -- (-1.233,1.927) -- (-1.240,1.919) -- (-1.248,1.911) -- (-1.256,1.903) -- (-1.263,1.895) -- (-1.271,1.886) -- (-1.279,1.878) -- (-1.287,1.869) -- (-1.294,1.860) -- (-1.302,1.851) -- (-1.310,1.841) -- (-1.318,1.832) -- (-1.326,1.822) -- (-1.334,1.812) -- (-1.342,1.802) -- (-1.350,1.792) -- (-1.358,1.782) -- (-1.367,1.771) -- (-1.375,1.760) -- (-1.383,1.749) -- (-1.391,1.738) -- (-1.399,1.726) -- (-1.407,1.714) -- (-1.413,1.706) -- (-1.418,1.698) -- (-1.424,1.690) -- (-1.429,1.682) -- (-1.435,1.674) -- (-1.440,1.665) -- (-1.446,1.657) -- (-1.451,1.648) -- (-1.457,1.639) -- (-1.462,1.630) -- (-1.468,1.621) -- (-1.473,1.612) -- (-1.479,1.603) -- (-1.484,1.593) -- (-1.490,1.584) -- (-1.495,1.574) -- (-1.501,1.565) -- (-1.506,1.555) -- (-1.512,1.545) -- (-1.517,1.535) -- (-1.522,1.524) -- (-1.528,1.514) -- (-1.533,1.504) -- (-1.538,1.493) -- (-1.544,1.482) -- (-1.549,1.471) -- (-1.554,1.461) -- (-1.560,1.449) -- (-1.565,1.438) -- (-1.570,1.427) -- (-1.575,1.415) -- (-1.581,1.404) -- (-1.586,1.392) -- (-1.591,1.380) -- (-1.596,1.368) -- (-1.601,1.356) -- (-1.606,1.343) -- (-1.611,1.331) -- (-1.616,1.318) -- (-1.621,1.306) -- (-1.626,1.293) -- (-1.630,1.280) -- (-1.635,1.267) -- (-1.640,1.253) -- (-1.645,1.240) -- (-1.649,1.226) -- (-1.654,1.213) -- (-1.658,1.199) -- (-1.663,1.185) -- (-1.667,1.171) -- (-1.671,1.157) -- (-1.676,1.142) -- (-1.680,1.128) -- (-1.684,1.113) -- (-1.688,1.098) -- (-1.692,1.083) -- (-1.696,1.068) -- (-1.700,1.053) -- (-1.704,1.037) -- (-1.707,1.022) -- (-1.711,1.006) -- (-1.714,0.991) -- (-1.718,0.975) -- (-1.721,0.959) -- (-1.725,0.943) -- (-1.728,0.926) -- (-1.731,0.910) -- (-1.734,0.893) -- (-1.737,0.877) -- (-1.740,0.860) -- (-1.742,0.843) -- (-1.745,0.826) -- (-1.747,0.809) -- (-1.750,0.792) -- (-1.752,0.774) -- (-1.754,0.757) -- (-1.756,0.739) -- (-1.758,0.722) -- (-1.760,0.704) -- (-1.762,0.686) -- (-1.763,0.668) -- (-1.765,0.650) -- (-1.766,0.632) -- (-1.768,0.613) -- (-1.769,0.595) -- (-1.770,0.577) -- (-1.771,0.558) -- (-1.772,0.539) -- (-1.772,0.521) -- (-1.773,0.502) -- (-1.773,0.483) -- (-1.773,0.464) -- (-1.774,0.446) -- (-1.774,0.427) -- (-1.774,0.408) -- (-1.773,0.388) -- (-1.773,0.369) -- (-1.772,0.350) -- (-1.772,0.341) -- (-1.772,0.331) -- (-1.772,0.321) -- (-1.771,0.312) -- (-1.771,0.302) -- (-1.770,0.292) -- (-1.770,0.283) -- (-1.769,0.273) -- (-1.769,0.263) -- (-1.768,0.254) -- (-1.767,0.244) -- (-1.767,0.234) -- (-1.766,0.225) -- (-1.765,0.215) -- (-1.765,0.205) -- (-1.764,0.196) -- (-1.763,0.186) -- (-1.762,0.176) -- (-1.761,0.167) -- (-1.760,0.157) -- (-1.759,0.147) -- (-1.758,0.138) -- (-1.757,0.128) -- (-1.756,0.118) -- (-1.755,0.108) -- (-1.754,0.099) -- (-1.753,0.089) -- (-1.752,0.079) -- (-1.751,0.070) -- (-1.750,0.060) -- (-1.748,0.050) -- (-1.747,0.041) -- (-1.746,0.031) -- (-1.744,0.021) -- (-1.743,0.012) -- (-1.741,0.002) -- (-1.740,-0.007) -- (-1.739,-0.017) -- (-1.737,-0.027) -- (-1.735,-0.036) -- (-1.734,-0.046) -- (-1.732,-0.055) -- (-1.731,-0.065) -- (-1.729,-0.075) -- (-1.727,-0.084) -- (-1.726,-0.094) -- (-1.724,-0.103) -- (-1.722,-0.113) -- (-1.720,-0.122) -- (-1.718,-0.132) -- (-1.717,-0.141) -- (-1.715,-0.151) -- (-1.713,-0.160) -- (-1.711,-0.169) -- (-1.709,-0.179) -- (-1.707,-0.188) -- (-1.703,-0.207) -- (-1.699,-0.226) -- (-1.694,-0.244) -- (-1.690,-0.263) -- (-1.685,-0.281) -- (-1.681,-0.300) -- (-1.676,-0.318) -- (-1.671,-0.336) -- (-1.666,-0.354) -- (-1.661,-0.372) -- (-1.656,-0.390) -- (-1.651,-0.408) -- (-1.646,-0.426) -- (-1.640,-0.443) -- (-1.635,-0.461) -- (-1.629,-0.478) -- (-1.624,-0.495) -- (-1.618,-0.513) -- (-1.612,-0.530) -- (-1.607,-0.547) -- (-1.601,-0.564) -- (-1.595,-0.580) -- (-1.589,-0.597) -- (-1.583,-0.613) -- (-1.577,-0.630) -- (-1.570,-0.646) -- (-1.564,-0.662) -- (-1.558,-0.678) -- (-1.551,-0.694) -- (-1.545,-0.710) -- (-1.538,-0.725) -- (-1.532,-0.741) -- (-1.525,-0.756) -- (-1.519,-0.771) -- (-1.512,-0.786) -- (-1.505,-0.801) -- (-1.499,-0.816) -- (-1.492,-0.831) -- (-1.485,-0.846) -- (-1.478,-0.860) -- (-1.471,-0.874) -- (-1.464,-0.888) -- (-1.457,-0.902) -- (-1.451,-0.916) -- (-1.444,-0.930) -- (-1.437,-0.944) -- (-1.430,-0.957) -- (-1.422,-0.971) -- (-1.415,-0.984) -- (-1.408,-0.997) -- (-1.401,-1.010) -- (-1.394,-1.023) -- (-1.387,-1.035) -- (-1.380,-1.048) -- (-1.373,-1.060) -- (-1.366,-1.073) -- (-1.358,-1.085) -- (-1.351,-1.097) -- (-1.344,-1.109) -- (-1.337,-1.121) -- (-1.330,-1.132) -- (-1.323,-1.144) -- (-1.315,-1.155) -- (-1.308,-1.167) -- (-1.301,-1.178) -- (-1.294,-1.189) -- (-1.287,-1.200) -- (-1.280,-1.210) -- (-1.272,-1.221) -- (-1.265,-1.232) -- (-1.258,-1.242) -- (-1.251,-1.252) -- (-1.244,-1.263) -- (-1.237,-1.273) -- (-1.230,-1.283) -- (-1.223,-1.292) -- (-1.216,-1.302) -- (-1.209,-1.312) -- (-1.202,-1.321) -- (-1.195,-1.331) -- (-1.188,-1.340) -- (-1.181,-1.349) -- (-1.174,-1.358) -- (-1.167,-1.367) -- (-1.160,-1.376) -- (-1.153,-1.385) -- (-1.146,-1.394) -- (-1.139,-1.402) -- (-1.132,-1.411) -- (-1.125,-1.419) -- (-1.118,-1.427) -- (-1.112,-1.435) -- (-1.105,-1.443) -- (-1.098,-1.451) -- (-1.091,-1.459) -- (-1.085,-1.467) -- (-1.078,-1.475) -- (-1.071,-1.482) -- (-1.065,-1.490) -- (-1.058,-1.497) -- (-1.051,-1.504) -- (-1.045,-1.512) -- (-1.038,-1.519) -- (-1.032,-1.526) -- (-1.022,-1.536) -- (-1.012,-1.547) -- (-1.003,-1.557) -- (-0.993,-1.567) -- (-0.984,-1.576) -- (-0.974,-1.586) -- (-0.965,-1.595) -- (-0.956,-1.604) -- (-0.946,-1.613) -- (-0.937,-1.622) -- (-0.928,-1.631) -- (-0.919,-1.639) -- (-0.910,-1.648) -- (-0.901,-1.656) -- (-0.892,-1.664) -- (-0.883,-1.672) -- (-0.874,-1.680) -- (-0.865,-1.688) -- (-0.857,-1.695) -- (-0.848,-1.702) -- (-0.840,-1.710) -- (-0.831,-1.717) -- (-0.823,-1.724) -- (-0.814,-1.731) -- (-0.806,-1.737) -- (-0.797,-1.744) -- (-0.789,-1.751) -- (-0.781,-1.757) -- (-0.773,-1.763) -- (-0.765,-1.769) -- (-0.757,-1.775) -- (-0.749,-1.781) -- (-0.741,-1.787) -- (-0.733,-1.793) -- (-0.725,-1.799) -- (-0.715,-1.806) -- (-0.705,-1.813) -- (-0.694,-1.820) -- (-0.684,-1.827) -- (-0.674,-1.834) -- (-0.664,-1.840) -- (-0.655,-1.847) -- (-0.645,-1.853) -- (-0.635,-1.859) -- (-0.626,-1.865) -- (-0.616,-1.871) -- (-0.607,-1.876) -- (-0.597,-1.882) -- (-0.588,-1.887) -- (-0.579,-1.892) -- (-0.570,-1.898) -- (-0.561,-1.903) -- (-0.552,-1.908) -- (-0.543,-1.912) -- (-0.534,-1.917) -- (-0.525,-1.922) -- (-0.517,-1.926) -- (-0.508,-1.930) -- (-0.497,-1.936) -- (-0.487,-1.941) -- (-0.476,-1.946) -- (-0.466,-1.951) -- (-0.455,-1.955) -- (-0.445,-1.960) -- (-0.435,-1.964) -- (-0.425,-1.969) -- (-0.415,-1.973) -- (-0.405,-1.977) -- (-0.395,-1.981) -- (-0.385,-1.984) -- (-0.376,-1.988) -- (-0.366,-1.991) -- (-0.356,-1.995) -- (-0.347,-1.998) -- (-0.337,-2.001) -- (-0.328,-2.004) -- (-0.319,-2.007) -- (-0.310,-2.010) -- (-0.299,-2.014) -- (-0.288,-2.017) -- (-0.277,-2.020) -- (-0.266,-2.023) -- (-0.255,-2.026) -- (-0.245,-2.028) -- (-0.234,-2.031) -- (-0.223,-2.033) -- (-0.213,-2.035) -- (-0.203,-2.038) -- (-0.192,-2.040) -- (-0.182,-2.041) -- (-0.172,-2.043) -- (-0.161,-2.045) -- (-0.151,-2.047) -- (-0.141,-2.048) -- (-0.131,-2.049) -- (-0.121,-2.051) -- (-0.111,-2.052) -- (-0.101,-2.053) -- (-0.091,-2.054) -- (-0.081,-2.055) -- (-0.071,-2.055) -- (-0.061,-2.056) -- (-0.051,-2.057) -- (-0.041,-2.057) -- (-0.031,-2.057) -- (-0.021,-2.058) -- (-0.012,-2.058) -- (-0.002,-2.058) -- (0.008,-2.058) -- (0.018,-2.058) -- (0.028,-2.058) -- (0.038,-2.057) -- (0.048,-2.057) -- (0.058,-2.056) -- (0.068,-2.056) -- (0.078,-2.055) -- (0.088,-2.054) -- (0.097,-2.053) -- (0.108,-2.052) -- (0.118,-2.051) -- (0.128,-2.050) -- (0.138,-2.049) -- (0.148,-2.047) -- (0.158,-2.046) -- (0.168,-2.044) -- (0.178,-2.042) -- (0.189,-2.040) -- (0.199,-2.038) -- (0.210,-2.036) -- (0.220,-2.034) -- (0.230,-2.031) -- (0.241,-2.029) -- (0.252,-2.026) -- (0.262,-2.024) -- (0.273,-2.021) -- (0.284,-2.018) -- (0.295,-2.015) -- (0.306,-2.011) -- (0.315,-2.009) -- (0.324,-2.006) -- (0.334,-2.003) -- (0.343,-1.999) -- (0.353,-1.996) -- (0.362,-1.993) -- (0.372,-1.989) -- (0.381,-1.986) -- (0.391,-1.982) -- (0.401,-1.978) -- (0.411,-1.974) -- (0.421,-1.970) -- (0.431,-1.966) -- (0.441,-1.962) -- (0.451,-1.957) -- (0.462,-1.953) -- (0.472,-1.948) -- (0.482,-1.943) -- (0.493,-1.938) -- (0.504,-1.933) -- (0.512,-1.928) -- (0.521,-1.924) -- (0.530,-1.919) -- (0.538,-1.915) -- (0.547,-1.910) -- (0.556,-1.905) -- (0.565,-1.900) -- (0.574,-1.895) -- (0.584,-1.890) -- (0.593,-1.885) -- (0.602,-1.879) -- (0.611,-1.873) -- (0.621,-1.868) -- (0.630,-1.862) -- (0.640,-1.856) -- (0.650,-1.850) -- (0.660,-1.843) -- (0.669,-1.837) -- (0.679,-1.830) -- (0.689,-1.824) -- (0.699,-1.817) -- (0.710,-1.810) -- (0.720,-1.802) -- (0.730,-1.795) -- (0.738,-1.789) -- (0.746,-1.783) -- (0.754,-1.777) -- (0.762,-1.771) -- (0.770,-1.765) -- (0.778,-1.759) -- (0.786,-1.753) -- (0.795,-1.746) -- (0.803,-1.740) -- (0.811,-1.733) -- (0.820,-1.726) -- (0.828,-1.719) -- (0.837,-1.712) -- (0.845,-1.705) -- (0.854,-1.698) -- (0.863,-1.690) -- (0.871,-1.682) -- (0.880,-1.675) -- (0.889,-1.667) -- (0.898,-1.659) -- (0.907,-1.651) -- (0.916,-1.642) -- (0.925,-1.634) -- (0.934,-1.625) -- (0.943,-1.616) -- (0.952,-1.607) -- (0.962,-1.598) -- (0.971,-1.589) -- (0.981,-1.579) -- (0.990,-1.570) -- (1.000,-1.560) -- (1.009,-1.550) -- (1.019,-1.540) -- (1.029,-1.529) -- (1.038,-1.519) -- (1.045,-1.512) -- (1.051,-1.504) -- (1.058,-1.497) -- (1.065,-1.490) -- (1.071,-1.482) -- (1.078,-1.475) -- (1.085,-1.467) -- (1.091,-1.459) -- (1.098,-1.451) -- (1.105,-1.443) -- (1.112,-1.435) -- (1.118,-1.427) -- (1.125,-1.419) -- (1.132,-1.411) -- (1.139,-1.402) -- (1.146,-1.394) -- (1.153,-1.385) -- (1.160,-1.376) -- (1.167,-1.367) -- (1.174,-1.358) -- (1.181,-1.349) -- (1.188,-1.340) -- (1.195,-1.331) -- (1.202,-1.321) -- (1.209,-1.312) -- (1.216,-1.302) -- (1.223,-1.292) -- (1.230,-1.283) -- (1.237,-1.273) -- (1.244,-1.263) -- (1.251,-1.252) -- (1.258,-1.242) -- (1.265,-1.232) -- (1.272,-1.221) -- (1.280,-1.210) -- (1.287,-1.200) -- (1.294,-1.189) -- (1.301,-1.178) -- (1.308,-1.167) -- (1.315,-1.155) -- (1.323,-1.144) -- (1.330,-1.132) -- (1.337,-1.121) -- (1.344,-1.109) -- (1.351,-1.097) -- (1.358,-1.085) -- (1.366,-1.073) -- (1.373,-1.060) -- (1.380,-1.048) -- (1.387,-1.035) -- (1.394,-1.023) -- (1.401,-1.010) -- (1.408,-0.997) -- (1.415,-0.984) -- (1.422,-0.971) -- (1.430,-0.957) -- (1.437,-0.944) -- (1.444,-0.930) -- (1.451,-0.916) -- (1.457,-0.902) -- (1.464,-0.888) -- (1.471,-0.874) -- (1.478,-0.860) -- (1.485,-0.846) -- (1.492,-0.831) -- (1.499,-0.816) -- (1.505,-0.801) -- (1.512,-0.786) -- (1.519,-0.771) -- (1.525,-0.756) -- (1.532,-0.741) -- (1.538,-0.725) -- (1.545,-0.710) -- (1.551,-0.694) -- (1.558,-0.678) -- (1.564,-0.662) -- (1.570,-0.646) -- (1.577,-0.630) -- (1.583,-0.613) -- (1.589,-0.597) -- (1.595,-0.580) -- (1.601,-0.564) -- (1.607,-0.547) -- (1.612,-0.530) -- (1.618,-0.513) -- (1.624,-0.495) -- (1.629,-0.478) -- (1.635,-0.461) -- (1.640,-0.443) -- (1.646,-0.426) -- (1.651,-0.408) -- (1.656,-0.390) -- (1.661,-0.372) -- (1.666,-0.354) -- (1.671,-0.336) -- (1.676,-0.318) -- (1.681,-0.300) -- (1.685,-0.281) -- (1.690,-0.263) -- (1.694,-0.244) -- (1.699,-0.226) -- (1.703,-0.207) -- (1.707,-0.188) -- (1.709,-0.179) -- (1.711,-0.169) -- (1.713,-0.160) -- (1.715,-0.151) -- (1.717,-0.141) -- (1.718,-0.132) -- (1.720,-0.122) -- (1.722,-0.113) -- (1.724,-0.103) -- (1.726,-0.094) -- (1.727,-0.084) -- (1.729,-0.075) -- (1.731,-0.065) -- (1.732,-0.055) -- (1.734,-0.046) -- (1.735,-0.036) -- (1.737,-0.027) -- (1.739,-0.017) -- (1.740,-0.007) -- (1.741,0.002) -- (1.743,0.012) -- (1.744,0.021) -- (1.746,0.031) -- (1.747,0.041) -- (1.748,0.050) -- (1.750,0.060) -- (1.751,0.070) -- (1.752,0.079) -- (1.753,0.089) -- (1.754,0.099) -- (1.755,0.108) -- (1.756,0.118) -- (1.757,0.128) -- (1.758,0.138) -- (1.759,0.147) -- (1.760,0.157) -- (1.761,0.167) -- (1.762,0.176) -- (1.763,0.186) -- (1.764,0.196) -- (1.765,0.205) -- (1.765,0.215) -- (1.766,0.225) -- (1.767,0.234) -- (1.767,0.244) -- (1.768,0.254) -- (1.769,0.263) -- (1.769,0.273) -- (1.770,0.283) -- (1.770,0.292) -- (1.771,0.302) -- (1.771,0.312) -- (1.772,0.321) -- (1.772,0.331) -- (1.772,0.341) -- (1.772,0.350) -- (1.773,0.369) -- (1.773,0.388) -- (1.774,0.408) -- (1.774,0.427) -- (1.774,0.446) -- (1.773,0.464) -- (1.773,0.483) -- (1.773,0.502) -- (1.772,0.521) -- (1.772,0.539) -- (1.771,0.558) -- (1.770,0.577) -- (1.769,0.595) -- (1.768,0.613) -- (1.766,0.632) -- (1.765,0.650) -- (1.763,0.668) -- (1.762,0.686) -- (1.760,0.704) -- (1.758,0.722) -- (1.756,0.739) -- (1.754,0.757) -- (1.752,0.774) -- (1.750,0.792) -- (1.747,0.809) -- (1.745,0.826) -- (1.742,0.843) -- (1.740,0.860) -- (1.737,0.877) -- (1.734,0.893) -- (1.731,0.910) -- (1.728,0.926) -- (1.725,0.943) -- (1.721,0.959) -- (1.718,0.975) -- (1.714,0.991) -- (1.711,1.006) -- (1.707,1.022) -- (1.704,1.037) -- (1.700,1.053) -- (1.696,1.068) -- (1.692,1.083) -- (1.688,1.098) -- (1.684,1.113) -- (1.680,1.128) -- (1.676,1.142) -- (1.671,1.157) -- (1.667,1.171) -- (1.663,1.185) -- (1.658,1.199) -- (1.654,1.213) -- (1.649,1.226) -- (1.645,1.240) -- (1.640,1.253) -- (1.635,1.267) -- (1.630,1.280) -- (1.626,1.293) -- (1.621,1.306) -- (1.616,1.318) -- (1.611,1.331) -- (1.606,1.343) -- (1.601,1.356) -- (1.596,1.368) -- (1.591,1.380) -- (1.586,1.392) -- (1.581,1.404) -- (1.575,1.415) -- (1.570,1.427) -- (1.565,1.438) -- (1.560,1.449) -- (1.554,1.461) -- (1.549,1.471) -- (1.544,1.482) -- (1.538,1.493) -- (1.533,1.504) -- (1.528,1.514) -- (1.522,1.524) -- (1.517,1.535) -- (1.512,1.545) -- (1.506,1.555) -- (1.501,1.565) -- (1.495,1.574) -- (1.490,1.584) -- (1.484,1.593) -- (1.479,1.603) -- (1.473,1.612) -- (1.468,1.621) -- (1.462,1.630) -- (1.457,1.639) -- (1.451,1.648) -- (1.446,1.657) -- (1.440,1.665) -- (1.435,1.674) -- (1.429,1.682) -- (1.424,1.690) -- (1.418,1.698) -- (1.413,1.706) -- (1.407,1.714) -- (1.399,1.726) -- (1.391,1.738) -- (1.383,1.749) -- (1.375,1.760) -- (1.367,1.771) -- (1.358,1.782) -- (1.350,1.792) -- (1.342,1.802) -- (1.334,1.812) -- (1.326,1.822) -- (1.318,1.832) -- (1.310,1.841) -- (1.302,1.851) -- (1.294,1.860) -- (1.287,1.869) -- (1.279,1.878) -- (1.271,1.886) -- (1.263,1.895) -- (1.256,1.903) -- (1.248,1.911) -- (1.240,1.919) -- (1.233,1.927) -- (1.225,1.934) -- (1.218,1.942) -- (1.210,1.949) -- (1.203,1.957) -- (1.196,1.964) -- (1.188,1.971) -- (1.181,1.977) -- (1.174,1.984) -- (1.167,1.991) -- (1.159,1.997) -- (1.150,2.006) -- (1.141,2.014) -- (1.131,2.022) -- (1.122,2.030) -- (1.113,2.037) -- (1.104,2.045) -- (1.095,2.052) -- (1.086,2.059) -- (1.077,2.066) -- (1.069,2.073) -- (1.060,2.079) -- (1.052,2.086) -- (1.043,2.092) -- (1.035,2.098) -- (1.026,2.104) -- (1.018,2.110) -- (1.010,2.115) -- (1.002,2.121) -- (0.994,2.126) -- (0.984,2.133) -- (0.974,2.139) -- (0.965,2.146) -- (0.955,2.152) -- (0.946,2.157) -- (0.937,2.163) -- (0.927,2.169) -- (0.918,2.174) -- (0.909,2.179) -- (0.901,2.184) -- (0.892,2.189) -- (0.883,2.194) -- (0.875,2.199) -- (0.865,2.204) -- (0.855,2.210) -- (0.845,2.215) -- (0.835,2.220) -- (0.826,2.224) -- (0.816,2.229) -- (0.807,2.233) -- (0.798,2.238) -- (0.789,2.242) -- (0.780,2.246) -- (0.771,2.250) -- (0.761,2.255) -- (0.751,2.259) -- (0.741,2.263) -- (0.732,2.267) -- (0.722,2.271) -- (0.713,2.275) -- (0.704,2.278) -- (0.695,2.282) -- (0.685,2.286) -- (0.675,2.289) -- (0.665,2.293) -- (0.655,2.296) -- (0.646,2.300) -- (0.637,2.303) -- (0.628,2.306) -- (0.617,2.309) -- (0.607,2.312) -- (0.598,2.316) -- (0.588,2.319) -- (0.579,2.321) -- (0.570,2.324) -- (0.559,2.327) -- (0.550,2.330) -- (0.540,2.332) -- (0.530,2.335) -- (0.521,2.337) -- (0.511,2.340) -- (0.501,2.343) -- (0.492,2.345) -- (0.482,2.347) -- (0.472,2.350) -- (0.462,2.352) -- (0.452,2.354) -- (0.443,2.356) -- (0.433,2.358) -- (0.423,2.360) -- (0.413,2.362) -- (0.404,2.364) -- (0.394,2.366) -- (0.384,2.367) -- (0.375,2.369) -- (0.365,2.371) -- (0.355,2.372) -- (0.345,2.374) -- (0.335,2.375) -- (0.326,2.377) -- (0.316,2.378) -- (0.306,2.380) -- (0.296,2.381) -- (0.286,2.382) -- (0.277,2.384) -- (0.267,2.385) -- (0.257,2.386) -- (0.247,2.387) -- (0.238,2.388) -- (0.228,2.389) -- (0.218,2.390) -- (0.209,2.391) -- (0.199,2.392) -- (0.189,2.392) -- (0.179,2.393) -- (0.169,2.394) -- (0.160,2.395) -- (0.150,2.395) -- (0.140,2.396) -- (0.130,2.396) -- (0.121,2.397) -- (0.111,2.397) -- (0.101,2.398) -- (0.091,2.398) -- (0.082,2.399) -- (0.072,2.399) -- (0.062,2.399) -- (0.053,2.399) -- (0.043,2.400) -- (0.033,2.400) -- (0.024,2.400) -- (0.014,2.400) -- (0.004,2.400) -- (0.000,2.400);}
\newcommand{\figOneSectionZero}{\draw[thick,gray] (0.000,-2.160) -- (0.000,-1.884);}
\newcommand{\figOneCentralSection}{\draw[thick,gray] (0.000,2.058) -- (0.000,2.400);}
\newcommand{\figOneVertex}{(0.000,-2.058)}
\newcommand{\figOneChartOrbits}[1]{\draw[#1] (8.852,1.888) -- (8.830,1.864) -- (8.808,1.841) -- (8.787,1.818) -- (8.766,1.797) -- (8.746,1.775) -- (8.727,1.754) -- (8.707,1.734) -- (8.689,1.715) -- (8.670,1.695) -- (8.653,1.677) -- (8.635,1.658) -- (8.618,1.640) -- (8.593,1.614) -- (8.569,1.589) -- (8.546,1.565) -- (8.523,1.542) -- (8.501,1.519) -- (8.480,1.497) -- (8.459,1.476) -- (8.439,1.455) -- (8.419,1.435) -- (8.400,1.415) -- (8.382,1.397) -- (8.363,1.378) -- (8.346,1.360) -- (8.329,1.343) -- (8.306,1.321) -- (8.285,1.299) -- (8.264,1.278) -- (8.244,1.258) -- (8.224,1.239) -- (8.205,1.220) -- (8.186,1.202) -- (8.168,1.184) -- (8.151,1.167) -- (8.134,1.150) -- (8.113,1.130) -- (8.093,1.111) -- (8.074,1.092) -- (8.055,1.074) -- (8.036,1.056) -- (8.019,1.040) -- (8.002,1.023) -- (7.982,1.004) -- (7.962,0.986) -- (7.943,0.969) -- (7.925,0.952) -- (7.908,0.936) -- (7.888,0.917) -- (7.869,0.900) -- (7.850,0.883) -- (7.832,0.867) -- (7.815,0.851) -- (7.796,0.834) -- (7.777,0.817) -- (7.759,0.801) -- (7.742,0.786) -- (7.723,0.769) -- (7.705,0.754) -- (7.687,0.738) -- (7.669,0.722) -- (7.651,0.707) -- (7.633,0.692) -- (7.614,0.676) -- (7.597,0.661) -- (7.579,0.647) -- (7.561,0.632) -- (7.544,0.618) -- (7.526,0.603) -- (7.508,0.589) -- (7.490,0.574) -- (7.473,0.560) -- (7.455,0.546) -- (7.437,0.533) -- (7.420,0.519) -- (7.401,0.505) -- (7.384,0.492) -- (7.366,0.479) -- (7.349,0.466) -- (7.331,0.452) -- (7.313,0.439) -- (7.295,0.427) -- (7.278,0.415) -- (7.260,0.402) -- (7.243,0.390) -- (7.226,0.378) -- (7.208,0.366) -- (7.190,0.355) -- (7.173,0.343) -- (7.155,0.331) -- (7.138,0.320) -- (7.120,0.309) -- (7.102,0.298) -- (7.085,0.288) -- (7.067,0.277) -- (7.050,0.267) -- (7.033,0.257) -- (7.015,0.247) -- (6.997,0.237) -- (6.980,0.227) -- (6.962,0.218) -- (6.944,0.209) -- (6.927,0.200) -- (6.909,0.191) -- (6.892,0.182) -- (6.874,0.174) -- (6.856,0.165) -- (6.839,0.157) -- (6.821,0.150) -- (6.804,0.142) -- (6.786,0.135) -- (6.769,0.128) -- (6.751,0.121) -- (6.734,0.114) -- (6.716,0.108) -- (6.698,0.102) -- (6.681,0.096) -- (6.663,0.090) -- (6.646,0.085) -- (6.628,0.080) -- (6.610,0.075) -- (6.593,0.070) -- (6.575,0.066) -- (6.558,0.062) -- (6.540,0.058) -- (6.522,0.054) -- (6.505,0.051) -- (6.487,0.048) -- (6.470,0.045) -- (6.452,0.042) -- (6.434,0.040) -- (6.417,0.038) -- (6.399,0.037) -- (6.382,0.035) -- (6.364,0.034) -- (6.346,0.033) -- (6.329,0.033) -- (6.311,0.032) -- (6.293,0.032) -- (6.276,0.033) -- (6.258,0.033) -- (6.241,0.034) -- (6.223,0.035) -- (6.205,0.036) -- (6.188,0.037) -- (6.170,0.039) -- (6.153,0.041) -- (6.135,0.043) -- (6.117,0.046) -- (6.100,0.048) -- (6.082,0.051) -- (6.064,0.054) -- (6.047,0.058) -- (6.029,0.061) -- (6.012,0.065) -- (5.994,0.069) -- (5.976,0.073) -- (5.959,0.078) -- (5.941,0.083) -- (5.924,0.088) -- (5.906,0.093) -- (5.888,0.098) -- (5.871,0.104) -- (5.853,0.110) -- (5.836,0.116) -- (5.818,0.122) -- (5.800,0.129) -- (5.783,0.135) -- (5.765,0.142) -- (5.748,0.149) -- (5.730,0.157) -- (5.713,0.164) -- (5.695,0.172) -- (5.678,0.180) -- (5.660,0.188) -- (5.643,0.197) -- (5.625,0.205) -- (5.607,0.214) -- (5.590,0.223) -- (5.572,0.232) -- (5.554,0.241) -- (5.537,0.251) -- (5.519,0.260) -- (5.502,0.270) -- (5.484,0.280) -- (5.466,0.290) -- (5.449,0.301) -- (5.431,0.311) -- (5.414,0.322) -- (5.396,0.333) -- (5.378,0.344) -- (5.361,0.355) -- (5.343,0.366) -- (5.326,0.377) -- (5.308,0.389) -- (5.290,0.401) -- (5.273,0.413) -- (5.255,0.425) -- (5.237,0.437) -- (5.220,0.449) -- (5.202,0.462) -- (5.184,0.474) -- (5.166,0.487) -- (5.148,0.500) -- (5.130,0.513) -- (5.113,0.526) -- (5.096,0.539) -- (5.078,0.552) -- (5.060,0.566) -- (5.043,0.579) -- (5.025,0.593) -- (5.007,0.607) -- (4.989,0.621) -- (4.971,0.635) -- (4.954,0.649) -- (4.936,0.663) -- (4.919,0.677) -- (4.901,0.691) -- (4.883,0.706) -- (4.865,0.721) -- (4.847,0.736) -- (4.829,0.750) -- (4.811,0.765) -- (4.794,0.780) -- (4.776,0.795) -- (4.758,0.810) -- (4.741,0.825) -- (4.723,0.840) -- (4.705,0.856) -- (4.686,0.872) -- (4.669,0.887) -- (4.651,0.903) -- (4.632,0.919) -- (4.615,0.935) -- (4.597,0.950) -- (4.579,0.967) -- (4.560,0.983) -- (4.543,0.999) -- (4.525,1.015) -- (4.507,1.031) -- (4.489,1.048) -- (4.470,1.065) -- (4.453,1.081) -- (4.435,1.097) -- (4.417,1.114) -- (4.398,1.131) -- (4.379,1.149) -- (4.360,1.167) -- (4.343,1.183) -- (4.325,1.200) -- (4.307,1.217) -- (4.288,1.234) -- (4.269,1.253) -- (4.250,1.271) -- (4.229,1.290) -- (4.212,1.307) -- (4.195,1.324) -- (4.177,1.341) -- (4.158,1.359) -- (4.139,1.377) -- (4.120,1.396) -- (4.100,1.415) -- (4.080,1.435) -- (4.059,1.455) -- (4.042,1.472) -- (4.024,1.489) -- (4.006,1.507) -- (3.988,1.525) -- (3.969,1.544) -- (3.950,1.563) -- (3.930,1.582) -- (3.910,1.602) -- (3.889,1.623) -- (3.868,1.644) -- (3.846,1.666) -- (3.824,1.688) -- (3.807,1.705) -- (3.790,1.723) -- (3.772,1.741) -- (3.754,1.759) -- (3.735,1.778) -- (3.716,1.797) -- (3.697,1.817) -- (3.677,1.837) -- (3.663,1.851); \draw[#1] (8.852,2.064) -- (8.828,2.037) -- (8.805,2.011) -- (8.783,1.985) -- (8.761,1.960) -- (8.740,1.936) -- (8.719,1.913) -- (8.699,1.890) -- (8.679,1.868) -- (8.660,1.846) -- (8.641,1.825) -- (8.623,1.805) -- (8.605,1.785) -- (8.587,1.765) -- (8.570,1.746) -- (8.545,1.719) -- (8.521,1.692) -- (8.498,1.666) -- (8.475,1.641) -- (8.453,1.617) -- (8.432,1.594) -- (8.411,1.571) -- (8.391,1.549) -- (8.371,1.528) -- (8.352,1.507) -- (8.333,1.487) -- (8.315,1.468) -- (8.298,1.449) -- (8.280,1.430) -- (8.258,1.407) -- (8.237,1.384) -- (8.216,1.362) -- (8.196,1.341) -- (8.176,1.320) -- (8.157,1.300) -- (8.138,1.281) -- (8.120,1.262) -- (8.103,1.244) -- (8.086,1.226) -- (8.065,1.205) -- (8.045,1.185) -- (8.026,1.165) -- (8.007,1.146) -- (7.989,1.127) -- (7.971,1.110) -- (7.954,1.093) -- (7.934,1.073) -- (7.915,1.053) -- (7.896,1.035) -- (7.878,1.017) -- (7.860,1.000) -- (7.843,0.984) -- (7.824,0.965) -- (7.806,0.947) -- (7.788,0.930) -- (7.770,0.913) -- (7.751,0.895) -- (7.732,0.878) -- (7.714,0.861) -- (7.697,0.845) -- (7.678,0.827) -- (7.660,0.810) -- (7.642,0.794) -- (7.624,0.777) -- (7.605,0.761) -- (7.588,0.745) -- (7.569,0.729) -- (7.551,0.713) -- (7.534,0.698) -- (7.516,0.682) -- (7.498,0.667) -- (7.480,0.651) -- (7.462,0.637) -- (7.444,0.621) -- (7.427,0.607) -- (7.409,0.592) -- (7.391,0.578) -- (7.373,0.563) -- (7.355,0.549) -- (7.338,0.535) -- (7.320,0.521) -- (7.302,0.507) -- (7.285,0.494) -- (7.267,0.481) -- (7.250,0.468) -- (7.232,0.455) -- (7.214,0.442) -- (7.197,0.429) -- (7.179,0.416) -- (7.161,0.404) -- (7.144,0.392) -- (7.126,0.380) -- (7.108,0.368) -- (7.091,0.357) -- (7.073,0.345) -- (7.056,0.334) -- (7.038,0.323) -- (7.020,0.312) -- (7.003,0.302) -- (6.986,0.291) -- (6.968,0.281) -- (6.950,0.271) -- (6.933,0.261) -- (6.915,0.252) -- (6.898,0.243) -- (6.880,0.233) -- (6.862,0.225) -- (6.845,0.216) -- (6.827,0.207) -- (6.810,0.199) -- (6.792,0.191) -- (6.774,0.184) -- (6.757,0.176) -- (6.739,0.169) -- (6.722,0.162) -- (6.704,0.156) -- (6.687,0.149) -- (6.669,0.143) -- (6.651,0.137) -- (6.634,0.132) -- (6.616,0.126) -- (6.598,0.121) -- (6.581,0.117) -- (6.563,0.112) -- (6.546,0.108) -- (6.528,0.104) -- (6.510,0.101) -- (6.493,0.097) -- (6.475,0.094) -- (6.458,0.092) -- (6.440,0.089) -- (6.422,0.087) -- (6.405,0.085) -- (6.387,0.084) -- (6.369,0.083) -- (6.352,0.082) -- (6.334,0.081) -- (6.316,0.081) -- (6.299,0.080) -- (6.281,0.081) -- (6.263,0.081) -- (6.246,0.082) -- (6.228,0.083) -- (6.211,0.084) -- (6.193,0.085) -- (6.175,0.087) -- (6.158,0.089) -- (6.140,0.091) -- (6.122,0.093) -- (6.105,0.096) -- (6.087,0.099) -- (6.069,0.102) -- (6.052,0.105) -- (6.034,0.109) -- (6.016,0.113) -- (5.999,0.117) -- (5.981,0.121) -- (5.963,0.125) -- (5.946,0.130) -- (5.928,0.135) -- (5.911,0.140) -- (5.893,0.145) -- (5.875,0.151) -- (5.858,0.157) -- (5.840,0.163) -- (5.823,0.169) -- (5.805,0.175) -- (5.787,0.182) -- (5.770,0.189) -- (5.752,0.196) -- (5.734,0.203) -- (5.717,0.210) -- (5.699,0.218) -- (5.682,0.226) -- (5.664,0.234) -- (5.646,0.242) -- (5.629,0.250) -- (5.611,0.259) -- (5.594,0.267) -- (5.576,0.276) -- (5.559,0.285) -- (5.541,0.295) -- (5.523,0.304) -- (5.506,0.314) -- (5.488,0.324) -- (5.470,0.334) -- (5.452,0.344) -- (5.435,0.354) -- (5.417,0.365) -- (5.399,0.376) -- (5.382,0.386) -- (5.364,0.397) -- (5.346,0.409) -- (5.329,0.420) -- (5.311,0.431) -- (5.293,0.443) -- (5.276,0.455) -- (5.258,0.466) -- (5.241,0.478) -- (5.223,0.490) -- (5.205,0.503) -- (5.187,0.515) -- (5.170,0.528) -- (5.152,0.540) -- (5.134,0.553) -- (5.117,0.566) -- (5.099,0.579) -- (5.081,0.593) -- (5.064,0.606) -- (5.046,0.620) -- (5.028,0.633) -- (5.010,0.647) -- (4.992,0.661) -- (4.974,0.675) -- (4.956,0.689) -- (4.938,0.704) -- (4.921,0.718) -- (4.903,0.732) -- (4.885,0.746) -- (4.868,0.761) -- (4.849,0.776) -- (4.832,0.790) -- (4.814,0.805) -- (4.795,0.821) -- (4.778,0.835) -- (4.760,0.851) -- (4.741,0.866) -- (4.724,0.881) -- (4.706,0.896) -- (4.688,0.912) -- (4.669,0.929) -- (4.652,0.944) -- (4.634,0.960) -- (4.616,0.976) -- (4.597,0.993) -- (4.580,1.008) -- (4.562,1.024) -- (4.544,1.040) -- (4.525,1.057) -- (4.506,1.075) -- (4.489,1.090) -- (4.471,1.106) -- (4.453,1.123) -- (4.434,1.140) -- (4.415,1.158) -- (4.395,1.176) -- (4.378,1.192) -- (4.360,1.209) -- (4.342,1.226) -- (4.323,1.244) -- (4.304,1.262) -- (4.284,1.280) -- (4.264,1.300) -- (4.247,1.316) -- (4.229,1.333) -- (4.211,1.350) -- (4.193,1.368) -- (4.174,1.386) -- (4.155,1.405) -- (4.135,1.424) -- (4.114,1.444) -- (4.093,1.465) -- (4.076,1.482) -- (4.058,1.499) -- (4.040,1.516) -- (4.022,1.535) -- (4.003,1.553) -- (3.984,1.572) -- (3.964,1.592) -- (3.944,1.612) -- (3.924,1.632) -- (3.902,1.653) -- (3.881,1.675) -- (3.858,1.698) -- (3.841,1.715) -- (3.824,1.732) -- (3.806,1.750) -- (3.788,1.769) -- (3.769,1.788) -- (3.750,1.807) -- (3.731,1.827) -- (3.711,1.847) -- (3.691,1.867) -- (3.670,1.889) -- (3.663,1.896); \draw[#1] (8.852,2.256) -- (8.827,2.225) -- (8.802,2.195) -- (8.778,2.166) -- (8.755,2.138) -- (8.733,2.111) -- (8.711,2.085) -- (8.689,2.059) -- (8.669,2.034) -- (8.648,2.010) -- (8.629,1.986) -- (8.609,1.963) -- (8.590,1.941) -- (8.572,1.919) -- (8.554,1.898) -- (8.537,1.877) -- (8.520,1.857) -- (8.503,1.837) -- (8.478,1.808) -- (8.455,1.781) -- (8.432,1.754) -- (8.410,1.728) -- (8.388,1.703) -- (8.367,1.679) -- (8.347,1.655) -- (8.327,1.633) -- (8.308,1.611) -- (8.289,1.589) -- (8.271,1.568) -- (8.253,1.548) -- (8.236,1.529) -- (8.219,1.509) -- (8.197,1.485) -- (8.176,1.461) -- (8.155,1.438) -- (8.135,1.416) -- (8.116,1.395) -- (8.097,1.374) -- (8.079,1.354) -- (8.061,1.335) -- (8.044,1.316) -- (8.023,1.293) -- (8.003,1.271) -- (7.984,1.250) -- (7.965,1.230) -- (7.946,1.210) -- (7.929,1.191) -- (7.911,1.173) -- (7.891,1.152) -- (7.872,1.131) -- (7.853,1.112) -- (7.835,1.093) -- (7.817,1.074) -- (7.800,1.057) -- (7.780,1.037) -- (7.762,1.018) -- (7.744,1.000) -- (7.726,0.982) -- (7.709,0.965) -- (7.690,0.947) -- (7.672,0.929) -- (7.655,0.911) -- (7.636,0.893) -- (7.617,0.875) -- (7.599,0.858) -- (7.582,0.842) -- (7.564,0.824) -- (7.546,0.808) -- (7.529,0.791) -- (7.510,0.775) -- (7.493,0.758) -- (7.474,0.742) -- (7.457,0.726) -- (7.439,0.710) -- (7.421,0.694) -- (7.404,0.679) -- (7.386,0.664) -- (7.369,0.649) -- (7.351,0.633) -- (7.334,0.619) -- (7.316,0.604) -- (7.298,0.589) -- (7.281,0.575) -- (7.263,0.561) -- (7.245,0.547) -- (7.228,0.533) -- (7.210,0.520) -- (7.193,0.506) -- (7.175,0.493) -- (7.157,0.480) -- (7.140,0.467) -- (7.122,0.454) -- (7.104,0.441) -- (7.087,0.429) -- (7.069,0.417) -- (7.052,0.405) -- (7.034,0.393) -- (7.016,0.382) -- (6.999,0.370) -- (6.981,0.359) -- (6.964,0.349) -- (6.946,0.338) -- (6.929,0.328) -- (6.911,0.318) -- (6.894,0.308) -- (6.876,0.298) -- (6.859,0.289) -- (6.841,0.280) -- (6.823,0.271) -- (6.806,0.262) -- (6.788,0.254) -- (6.770,0.246) -- (6.753,0.238) -- (6.735,0.230) -- (6.718,0.223) -- (6.700,0.216) -- (6.682,0.210) -- (6.665,0.203) -- (6.647,0.197) -- (6.630,0.191) -- (6.612,0.186) -- (6.594,0.181) -- (6.577,0.176) -- (6.559,0.171) -- (6.541,0.167) -- (6.524,0.163) -- (6.506,0.159) -- (6.488,0.156) -- (6.470,0.153) -- (6.453,0.150) -- (6.435,0.148) -- (6.417,0.146) -- (6.400,0.144) -- (6.382,0.142) -- (6.364,0.141) -- (6.347,0.140) -- (6.329,0.140) -- (6.311,0.139) -- (6.294,0.139) -- (6.276,0.139) -- (6.258,0.140) -- (6.241,0.141) -- (6.223,0.142) -- (6.205,0.143) -- (6.188,0.145) -- (6.170,0.147) -- (6.152,0.149) -- (6.135,0.151) -- (6.117,0.154) -- (6.099,0.156) -- (6.082,0.159) -- (6.064,0.163) -- (6.046,0.166) -- (6.029,0.170) -- (6.011,0.174) -- (5.993,0.178) -- (5.975,0.182) -- (5.958,0.187) -- (5.940,0.192) -- (5.922,0.197) -- (5.905,0.202) -- (5.887,0.208) -- (5.870,0.213) -- (5.852,0.219) -- (5.834,0.225) -- (5.817,0.232) -- (5.799,0.238) -- (5.781,0.245) -- (5.764,0.252) -- (5.746,0.259) -- (5.728,0.266) -- (5.711,0.274) -- (5.693,0.281) -- (5.675,0.289) -- (5.658,0.297) -- (5.640,0.306) -- (5.622,0.314) -- (5.604,0.323) -- (5.587,0.332) -- (5.569,0.341) -- (5.551,0.350) -- (5.533,0.359) -- (5.516,0.369) -- (5.498,0.379) -- (5.481,0.388) -- (5.463,0.398) -- (5.446,0.409) -- (5.428,0.419) -- (5.411,0.429) -- (5.393,0.440) -- (5.376,0.451) -- (5.358,0.462) -- (5.340,0.473) -- (5.322,0.485) -- (5.304,0.496) -- (5.287,0.508) -- (5.269,0.519) -- (5.251,0.531) -- (5.234,0.543) -- (5.216,0.556) -- (5.198,0.568) -- (5.181,0.580) -- (5.163,0.593) -- (5.146,0.605) -- (5.128,0.618) -- (5.110,0.631) -- (5.092,0.644) -- (5.075,0.658) -- (5.057,0.671) -- (5.040,0.684) -- (5.022,0.698) -- (5.003,0.712) -- (4.986,0.726) -- (4.968,0.740) -- (4.950,0.754) -- (4.932,0.768) -- (4.914,0.783) -- (4.896,0.797) -- (4.878,0.812) -- (4.860,0.827) -- (4.842,0.841) -- (4.825,0.856) -- (4.806,0.872) -- (4.789,0.886) -- (4.771,0.901) -- (4.753,0.917) -- (4.734,0.933) -- (4.717,0.948) -- (4.699,0.963) -- (4.681,0.979) -- (4.662,0.995) -- (4.643,1.012) -- (4.625,1.028) -- (4.607,1.044) -- (4.589,1.060) -- (4.570,1.077) -- (4.553,1.093) -- (4.535,1.109) -- (4.517,1.125) -- (4.499,1.142) -- (4.480,1.159) -- (4.460,1.177) -- (4.443,1.193) -- (4.426,1.209) -- (4.408,1.226) -- (4.390,1.243) -- (4.371,1.260) -- (4.351,1.278) -- (4.332,1.297) -- (4.311,1.316) -- (4.294,1.333) -- (4.276,1.350) -- (4.258,1.367) -- (4.239,1.385) -- (4.220,1.404) -- (4.201,1.423) -- (4.181,1.442) -- (4.160,1.462) -- (4.139,1.483) -- (4.122,1.500) -- (4.104,1.517) -- (4.086,1.535) -- (4.067,1.553) -- (4.048,1.572) -- (4.029,1.591) -- (4.009,1.610) -- (3.989,1.631) -- (3.968,1.651) -- (3.947,1.673) -- (3.925,1.695) -- (3.903,1.717) -- (3.885,1.734) -- (3.868,1.752) -- (3.850,1.770) -- (3.832,1.789) -- (3.813,1.808) -- (3.794,1.827) -- (3.774,1.847) -- (3.755,1.867) -- (3.734,1.888) -- (3.713,1.910) -- (3.692,1.932) -- (3.670,1.954) -- (3.663,1.962); \draw[#1] (8.852,0.560) -- (8.829,0.550) -- (8.807,0.541) -- (8.786,0.531) -- (8.765,0.522) -- (8.745,0.514) -- (8.725,0.505) -- (8.705,0.497) -- (8.687,0.489) -- (8.668,0.481) -- (8.650,0.474) -- (8.629,0.465) -- (8.608,0.456) -- (8.587,0.447) -- (8.567,0.439) -- (8.548,0.431) -- (8.529,0.423) -- (8.511,0.416) -- (8.493,0.409) -- (8.472,0.400) -- (8.452,0.392) -- (8.432,0.384) -- (8.413,0.377) -- (8.395,0.369) -- (8.377,0.362) -- (8.357,0.354) -- (8.337,0.347) -- (8.318,0.339) -- (8.299,0.332) -- (8.281,0.325) -- (8.261,0.318) -- (8.242,0.310) -- (8.224,0.303) -- (8.205,0.296) -- (8.188,0.290) -- (8.169,0.283) -- (8.150,0.276) -- (8.132,0.270) -- (8.113,0.263) -- (8.095,0.256) -- (8.077,0.250) -- (8.057,0.243) -- (8.039,0.236) -- (8.021,0.230) -- (8.002,0.224) -- (7.984,0.217) -- (7.965,0.211) -- (7.947,0.205) -- (7.929,0.199) -- (7.911,0.193) -- (7.893,0.187) -- (7.875,0.181) -- (7.858,0.176) -- (7.840,0.170) -- (7.821,0.164) -- (7.804,0.159) -- (7.786,0.153) -- (7.767,0.148) -- (7.750,0.143) -- (7.732,0.137) -- (7.714,0.132) -- (7.696,0.127) -- (7.678,0.122) -- (7.660,0.117) -- (7.642,0.112) -- (7.625,0.108) -- (7.607,0.103) -- (7.589,0.098) -- (7.571,0.094) -- (7.553,0.089) -- (7.535,0.085) -- (7.518,0.081) -- (7.500,0.077) -- (7.482,0.073) -- (7.464,0.069) -- (7.446,0.065) -- (7.429,0.061) -- (7.411,0.058) -- (7.393,0.054) -- (7.375,0.051) -- (7.358,0.047) -- (7.340,0.044) -- (7.322,0.041) -- (7.304,0.038) -- (7.287,0.035) -- (7.269,0.032) -- (7.251,0.030) -- (7.234,0.027) -- (7.216,0.025) -- (7.198,0.023) -- (7.181,0.020) -- (7.163,0.018) -- (7.145,0.017) -- (7.128,0.015) -- (7.110,0.013) -- (7.092,0.012) -- (7.075,0.010) -- (7.057,0.009) -- (7.039,0.008) -- (7.022,0.007) -- (7.004,0.006) -- (6.986,0.005) -- (6.969,0.004) -- (6.951,0.003) -- (6.934,0.003) -- (6.916,0.002) -- (6.898,0.002) -- (6.881,0.001) -- (6.863,0.001) -- (6.845,0.001) -- (6.828,0.001) -- (6.810,0.000) -- (6.793,0.000) -- (6.775,0.000) -- (6.757,0.000) -- (6.740,0.000) -- (6.722,0.000) -- (6.705,0.000) -- (6.687,0.000) -- (6.669,0.000) -- (6.652,0.000) -- (6.634,0.000) -- (6.617,0.000) -- (6.599,0.000) -- (6.581,0.000) -- (6.564,0.000) -- (6.546,0.000) -- (6.529,0.000) -- (6.511,0.000) -- (6.493,0.000) -- (6.476,0.000) -- (6.458,0.000) -- (6.441,0.000) -- (6.423,0.000) -- (6.405,0.000) -- (6.388,0.000) -- (6.378,0.000); \draw[#1] (8.852,1.120) -- (8.828,1.103) -- (8.805,1.086) -- (8.782,1.069) -- (8.760,1.053) -- (8.739,1.038) -- (8.718,1.023) -- (8.698,1.008) -- (8.678,0.994) -- (8.659,0.980) -- (8.640,0.967) -- (8.622,0.954) -- (8.604,0.941) -- (8.587,0.929) -- (8.565,0.913) -- (8.543,0.898) -- (8.521,0.883) -- (8.501,0.868) -- (8.481,0.854) -- (8.461,0.841) -- (8.443,0.828) -- (8.424,0.815) -- (8.406,0.803) -- (8.389,0.791) -- (8.368,0.776) -- (8.347,0.762) -- (8.327,0.749) -- (8.308,0.735) -- (8.289,0.723) -- (8.271,0.711) -- (8.254,0.699) -- (8.233,0.685) -- (8.213,0.672) -- (8.194,0.659) -- (8.175,0.647) -- (8.157,0.635) -- (8.140,0.623) -- (8.120,0.610) -- (8.100,0.598) -- (8.082,0.586) -- (8.064,0.574) -- (8.047,0.563) -- (8.027,0.551) -- (8.009,0.539) -- (7.991,0.528) -- (7.973,0.517) -- (7.954,0.505) -- (7.936,0.494) -- (7.918,0.483) -- (7.899,0.471) -- (7.881,0.460) -- (7.863,0.449) -- (7.845,0.438) -- (7.827,0.427) -- (7.809,0.417) -- (7.791,0.406) -- (7.774,0.396) -- (7.755,0.385) -- (7.738,0.375) -- (7.719,0.365) -- (7.702,0.355) -- (7.684,0.345) -- (7.665,0.334) -- (7.647,0.324) -- (7.629,0.314) -- (7.611,0.304) -- (7.593,0.295) -- (7.575,0.286) -- (7.558,0.276) -- (7.540,0.267) -- (7.522,0.257) -- (7.504,0.248) -- (7.486,0.239) -- (7.468,0.230) -- (7.450,0.222) -- (7.432,0.213) -- (7.415,0.204) -- (7.397,0.196) -- (7.379,0.188) -- (7.361,0.179) -- (7.344,0.171) -- (7.326,0.163) -- (7.308,0.155) -- (7.290,0.148) -- (7.272,0.140) -- (7.255,0.133) -- (7.237,0.126) -- (7.220,0.119) -- (7.202,0.112) -- (7.184,0.105) -- (7.166,0.098) -- (7.149,0.092) -- (7.131,0.086) -- (7.114,0.079) -- (7.096,0.074) -- (7.079,0.068) -- (7.061,0.062) -- (7.043,0.057) -- (7.026,0.052) -- (7.008,0.047) -- (6.990,0.042) -- (6.973,0.037) -- (6.955,0.033) -- (6.938,0.029) -- (6.920,0.025) -- (6.902,0.022) -- (6.885,0.018) -- (6.867,0.015) -- (6.850,0.013) -- (6.832,0.010) -- (6.814,0.008) -- (6.797,0.006) -- (6.779,0.005) -- (6.761,0.003) -- (6.744,0.002) -- (6.726,0.002) -- (6.709,0.001) -- (6.691,0.001) -- (6.673,0.000) -- (6.656,0.000) -- (6.638,0.000) -- (6.621,0.000) -- (6.603,0.000) -- (6.585,0.000) -- (6.568,0.000) -- (6.550,0.000) -- (6.533,0.000) -- (6.515,0.000) -- (6.497,0.000) -- (6.480,0.000) -- (6.462,0.000) -- (6.445,0.000) -- (6.427,0.000) -- (6.409,0.000) -- (6.392,0.000) -- (6.377,0.000); \draw[#1] (3.748,0.096) -- (3.768,0.097) -- (3.788,0.098) -- (3.807,0.098) -- (3.825,0.099) -- (3.843,0.100) -- (3.861,0.101) -- (3.881,0.102) -- (3.900,0.102) -- (3.918,0.103) -- (3.936,0.104) -- (3.956,0.105) -- (3.975,0.106) -- (3.994,0.108) -- (4.012,0.109) -- (4.031,0.110) -- (4.050,0.111) -- (4.068,0.112) -- (4.086,0.114) -- (4.104,0.115) -- (4.123,0.117) -- (4.140,0.118) -- (4.159,0.120) -- (4.177,0.121) -- (4.195,0.123) -- (4.213,0.125) -- (4.232,0.127) -- (4.250,0.129) -- (4.269,0.131) -- (4.287,0.134) -- (4.305,0.136) -- (4.323,0.139) -- (4.341,0.142) -- (4.359,0.145) -- (4.377,0.148) -- (4.395,0.151) -- (4.413,0.155) -- (4.431,0.159) -- (4.449,0.163) -- (4.467,0.168) -- (4.485,0.173) -- (4.503,0.179) -- (4.520,0.186) -- (4.538,0.193) -- (4.556,0.202) -- (4.573,0.212) -- (4.591,0.225) -- (4.608,0.242) -- (4.624,0.268) -- (4.631,0.326) -- (4.619,0.373) -- (4.604,0.407) -- (4.588,0.436) -- (4.572,0.462) -- (4.555,0.486) -- (4.539,0.509) -- (4.522,0.532) -- (4.505,0.553) -- (4.488,0.574) -- (4.471,0.594) -- (4.454,0.614) -- (4.437,0.635) -- (4.420,0.654) -- (4.403,0.673) -- (4.386,0.692) -- (4.368,0.712) -- (4.350,0.731) -- (4.333,0.750) -- (4.316,0.768) -- (4.298,0.787) -- (4.281,0.805) -- (4.264,0.824) -- (4.246,0.842) -- (4.229,0.860) -- (4.210,0.880) -- (4.192,0.898) -- (4.174,0.917) -- (4.157,0.935) -- (4.139,0.953) -- (4.121,0.972) -- (4.104,0.989) -- (4.086,1.007) -- (4.068,1.026) -- (4.049,1.045) -- (4.032,1.063) -- (4.015,1.080) -- (3.997,1.099) -- (3.978,1.118) -- (3.959,1.138) -- (3.942,1.155) -- (3.924,1.173) -- (3.906,1.191) -- (3.887,1.210) -- (3.868,1.230) -- (3.848,1.250) -- (3.831,1.267) -- (3.814,1.285) -- (3.796,1.303) -- (3.777,1.322) -- (3.758,1.341) -- (3.739,1.361) -- (3.718,1.381) -- (3.698,1.402) -- (3.681,1.419) -- (3.663,1.437); \draw[#1] (3.748,0.320) -- (3.766,0.325) -- (3.785,0.331) -- (3.803,0.337) -- (3.821,0.343) -- (3.839,0.350) -- (3.858,0.358) -- (3.876,0.367) -- (3.894,0.376) -- (3.912,0.387) -- (3.930,0.400) -- (3.948,0.415) -- (3.965,0.433) -- (3.981,0.457) -- (3.996,0.494) -- (4.000,0.556) -- (3.990,0.608) -- (3.976,0.649) -- (3.961,0.682) -- (3.945,0.713) -- (3.929,0.742) -- (3.912,0.769) -- (3.896,0.795) -- (3.878,0.820) -- (3.861,0.845) -- (3.843,0.870) -- (3.826,0.893) -- (3.808,0.916) -- (3.790,0.939) -- (3.773,0.961) -- (3.756,0.982) -- (3.738,1.004) -- (3.719,1.026) -- (3.702,1.047) -- (3.684,1.068) -- (3.665,1.090) -- (3.662,1.093); \draw[#1] (3.748,0.720) -- (3.734,0.758) -- (3.718,0.792) -- (3.702,0.823) -- (3.685,0.853) -- (3.669,0.880) -- (3.661,0.893);}
\newcommand{\figOneChartLeft}[1]{\draw[#1] (3.660,1.832) -- (3.673,1.818) -- (3.686,1.805) -- (3.700,1.791) -- (3.713,1.778) -- (3.726,1.765) -- (3.739,1.751) -- (3.752,1.738) -- (3.766,1.725) -- (3.779,1.711) -- (3.792,1.698) -- (3.805,1.685) -- (3.818,1.672) -- (3.832,1.658) -- (3.845,1.645) -- (3.858,1.632) -- (3.871,1.619) -- (3.884,1.606) -- (3.898,1.592) -- (3.911,1.579) -- (3.924,1.566) -- (3.937,1.553) -- (3.950,1.540) -- (3.964,1.527) -- (3.977,1.514) -- (3.990,1.501) -- (4.003,1.488) -- (4.016,1.475) -- (4.030,1.462) -- (4.043,1.449) -- (4.056,1.436) -- (4.069,1.423) -- (4.082,1.410) -- (4.096,1.398) -- (4.109,1.385) -- (4.122,1.372) -- (4.135,1.359) -- (4.148,1.346) -- (4.162,1.334) -- (4.175,1.321) -- (4.188,1.308) -- (4.201,1.296) -- (4.214,1.283) -- (4.228,1.270) -- (4.241,1.258) -- (4.254,1.245) -- (4.267,1.233) -- (4.280,1.220) -- (4.294,1.208) -- (4.307,1.195) -- (4.320,1.183) -- (4.333,1.170) -- (4.346,1.158) -- (4.360,1.146) -- (4.373,1.133) -- (4.386,1.121) -- (4.399,1.109) -- (4.412,1.097) -- (4.426,1.084) -- (4.439,1.072) -- (4.452,1.060) -- (4.465,1.048) -- (4.478,1.036) -- (4.492,1.024) -- (4.505,1.012) -- (4.518,1.000) -- (4.531,0.988) -- (4.544,0.976) -- (4.558,0.964) -- (4.571,0.952) -- (4.584,0.940) -- (4.597,0.929) -- (4.610,0.917) -- (4.624,0.905) -- (4.637,0.894) -- (4.650,0.882) -- (4.663,0.870) -- (4.676,0.859) -- (4.690,0.847) -- (4.703,0.836) -- (4.716,0.824) -- (4.729,0.813) -- (4.742,0.802) -- (4.756,0.790) -- (4.769,0.779) -- (4.782,0.768) -- (4.795,0.757) -- (4.808,0.746) -- (4.822,0.735) -- (4.835,0.724) -- (4.848,0.713) -- (4.861,0.702) -- (4.874,0.691) -- (4.888,0.680) -- (4.901,0.669) -- (4.914,0.658) -- (4.927,0.648) -- (4.940,0.637) -- (4.954,0.626) -- (4.967,0.616) -- (4.980,0.605) -- (4.993,0.595) -- (5.006,0.585) -- (5.020,0.574) -- (5.033,0.564) -- (5.046,0.554) -- (5.059,0.544) -- (5.072,0.534) -- (5.086,0.524) -- (5.099,0.514) -- (5.112,0.504) -- (5.125,0.494) -- (5.138,0.484) -- (5.152,0.475) -- (5.165,0.465) -- (5.178,0.455) -- (5.191,0.446) -- (5.204,0.436) -- (5.218,0.427) -- (5.231,0.418) -- (5.244,0.409) -- (5.257,0.399) -- (5.270,0.390) -- (5.284,0.381) -- (5.297,0.373) -- (5.310,0.364) -- (5.323,0.355) -- (5.336,0.346) -- (5.350,0.338) -- (5.363,0.329) -- (5.376,0.321) -- (5.389,0.312) -- (5.402,0.304) -- (5.416,0.296) -- (5.429,0.288) -- (5.442,0.280) -- (5.455,0.272) -- (5.468,0.264) -- (5.482,0.256) -- (5.495,0.249) -- (5.508,0.241) -- (5.521,0.234) -- (5.534,0.226) -- (5.548,0.219) -- (5.561,0.212) -- (5.574,0.205) -- (5.587,0.198) -- (5.600,0.191) -- (5.614,0.184) -- (5.627,0.177) -- (5.640,0.171) -- (5.653,0.164) -- (5.666,0.158) -- (5.680,0.152) -- (5.693,0.146) -- (5.706,0.140) -- (5.719,0.134) -- (5.732,0.128) -- (5.746,0.122) -- (5.759,0.117) -- (5.772,0.111) -- (5.785,0.106) -- (5.798,0.101) -- (5.812,0.096) -- (5.825,0.091) -- (5.838,0.086) -- (5.851,0.081) -- (5.864,0.077) -- (5.878,0.072) -- (5.891,0.068) -- (5.904,0.064) -- (5.917,0.059) -- (5.930,0.055) -- (5.944,0.052) -- (5.957,0.048) -- (5.970,0.044) -- (5.983,0.041) -- (5.996,0.038) -- (6.010,0.034) -- (6.023,0.031) -- (6.036,0.029) -- (6.049,0.026) -- (6.062,0.023) -- (6.076,0.021) -- (6.089,0.018) -- (6.102,0.016) -- (6.115,0.014) -- (6.128,0.012) -- (6.142,0.010) -- (6.155,0.009) -- (6.168,0.007) -- (6.181,0.006) -- (6.194,0.005) -- (6.208,0.004) -- (6.221,0.003) -- (6.234,0.002) -- (6.247,0.001) -- (6.260,0.001) -- (6.274,0.000) -- (6.287,0.000) -- (6.300,0.000);}
\newcommand{\figOneChartRight}[1]{\draw[#1] (6.300,0.000) -- (6.313,0.000) -- (6.326,0.000) -- (6.340,0.001) -- (6.353,0.001) -- (6.366,0.002) -- (6.379,0.003) -- (6.392,0.004) -- (6.406,0.005) -- (6.419,0.006) -- (6.432,0.007) -- (6.445,0.009) -- (6.458,0.010) -- (6.472,0.012) -- (6.485,0.014) -- (6.498,0.016) -- (6.511,0.018) -- (6.524,0.021) -- (6.538,0.023) -- (6.551,0.026) -- (6.564,0.029) -- (6.577,0.031) -- (6.590,0.034) -- (6.604,0.038) -- (6.617,0.041) -- (6.630,0.044) -- (6.643,0.048) -- (6.656,0.052) -- (6.670,0.055) -- (6.683,0.059) -- (6.696,0.064) -- (6.709,0.068) -- (6.722,0.072) -- (6.736,0.077) -- (6.749,0.081) -- (6.762,0.086) -- (6.775,0.091) -- (6.788,0.096) -- (6.802,0.101) -- (6.815,0.106) -- (6.828,0.111) -- (6.841,0.117) -- (6.854,0.122) -- (6.868,0.128) -- (6.881,0.134) -- (6.894,0.140) -- (6.907,0.146) -- (6.920,0.152) -- (6.934,0.158) -- (6.947,0.164) -- (6.960,0.171) -- (6.973,0.177) -- (6.986,0.184) -- (7.000,0.191) -- (7.013,0.198) -- (7.026,0.205) -- (7.039,0.212) -- (7.052,0.219) -- (7.066,0.226) -- (7.079,0.234) -- (7.092,0.241) -- (7.105,0.249) -- (7.118,0.256) -- (7.132,0.264) -- (7.145,0.272) -- (7.158,0.280) -- (7.171,0.288) -- (7.184,0.296) -- (7.198,0.304) -- (7.211,0.312) -- (7.224,0.321) -- (7.237,0.329) -- (7.250,0.338) -- (7.264,0.346) -- (7.277,0.355) -- (7.290,0.364) -- (7.303,0.373) -- (7.316,0.381) -- (7.330,0.390) -- (7.343,0.399) -- (7.356,0.409) -- (7.369,0.418) -- (7.382,0.427) -- (7.396,0.436) -- (7.409,0.446) -- (7.422,0.455) -- (7.435,0.465) -- (7.448,0.475) -- (7.462,0.484) -- (7.475,0.494) -- (7.488,0.504) -- (7.501,0.514) -- (7.514,0.524) -- (7.528,0.534) -- (7.541,0.544) -- (7.554,0.554) -- (7.567,0.564) -- (7.580,0.574) -- (7.594,0.585) -- (7.607,0.595) -- (7.620,0.605) -- (7.633,0.616) -- (7.646,0.626) -- (7.660,0.637) -- (7.673,0.648) -- (7.686,0.658) -- (7.699,0.669) -- (7.712,0.680) -- (7.726,0.691) -- (7.739,0.702) -- (7.752,0.713) -- (7.765,0.724) -- (7.778,0.735) -- (7.792,0.746) -- (7.805,0.757) -- (7.818,0.768) -- (7.831,0.779) -- (7.844,0.790) -- (7.858,0.802) -- (7.871,0.813) -- (7.884,0.824) -- (7.897,0.836) -- (7.910,0.847) -- (7.924,0.859) -- (7.937,0.870) -- (7.950,0.882) -- (7.963,0.894) -- (7.976,0.905) -- (7.990,0.917) -- (8.003,0.929) -- (8.016,0.940) -- (8.029,0.952) -- (8.042,0.964) -- (8.056,0.976) -- (8.069,0.988) -- (8.082,1.000) -- (8.095,1.012) -- (8.108,1.024) -- (8.122,1.036) -- (8.135,1.048) -- (8.148,1.060) -- (8.161,1.072) -- (8.174,1.084) -- (8.188,1.097) -- (8.201,1.109) -- (8.214,1.121) -- (8.227,1.133) -- (8.240,1.146) -- (8.254,1.158) -- (8.267,1.170) -- (8.280,1.183) -- (8.293,1.195) -- (8.306,1.208) -- (8.320,1.220) -- (8.333,1.233) -- (8.346,1.245) -- (8.359,1.258) -- (8.372,1.270) -- (8.386,1.283) -- (8.399,1.296) -- (8.412,1.308) -- (8.425,1.321) -- (8.438,1.334) -- (8.452,1.346) -- (8.465,1.359) -- (8.478,1.372) -- (8.491,1.385) -- (8.504,1.398) -- (8.518,1.410) -- (8.531,1.423) -- (8.544,1.436) -- (8.557,1.449) -- (8.570,1.462) -- (8.584,1.475) -- (8.597,1.488) -- (8.610,1.501) -- (8.623,1.514) -- (8.636,1.527) -- (8.650,1.540) -- (8.663,1.553) -- (8.676,1.566) -- (8.689,1.579) -- (8.702,1.592) -- (8.716,1.606) -- (8.729,1.619) -- (8.742,1.632) -- (8.755,1.645) -- (8.768,1.658) -- (8.782,1.672) -- (8.795,1.685) -- (8.808,1.698) -- (8.821,1.711) -- (8.834,1.725) -- (8.848,1.738) -- (8.861,1.751) -- (8.874,1.765) -- (8.887,1.778) -- (8.900,1.791) -- (8.914,1.805) -- (8.927,1.818) -- (8.940,1.832);}

\begin{tikzpicture}[>=Stealth,font=\small,
    mid/.style={postaction={decorate},decoration={markings,
      mark=at position #1 with {\arrow{Stealth}}}},
    midrev/.style={postaction={decorate},decoration={markings,
      mark=at position #1 with {\arrowreversed{Stealth}}}}]
  % (a) the Poincare disk
  \draw[gray!60] (0,0) circle (2.4);
  \figOneSpiral{gray,thin,mid=.12}
  \figOneGraphic{very thick,
    postaction={decorate},decoration={markings,
      mark=at position .3 with {\arrow{Stealth}},
      mark=at position .7 with {\arrow{Stealth}}}}
  \figOneSectionZero
  \figOneCentralSection
  \fill (0,0) circle (1.5pt) (0,2.4) circle (1.8pt);
  \node[above] at (0,2.4) {saddle--node};
  \node[right,xshift=4pt,yshift=-14pt] at \figOneVertex {$\ \Sigma_0$};
  \node[right] at (0,2.02) {$\Sigma_c$};
  \node[right] at (2.55,-.8) {$\Gamma(C_0)$};
  \node at (0,-2.85) {(a)};
  % (b) the projective chart near the saddle--node
  \draw[->,gray] (3.4,0)--(9.25,0) node[right] {$v$};
  \draw[->,gray] (6.3,0)--(6.3,2.75) node[above] {$z$};
  \figOneChartOrbits{gray,thin,mid=.55}
  \draw[very thick,mid=.55] (3.66,0)--(6.3,0);
  \draw[very thick,mid=.55] (8.94,0)--(6.3,0);
  \figOneChartLeft{very thick,midrev=.5}
  \figOneChartRight{very thick,midrev=.5}
  \draw[dashed] (6.3,0)--(6.3,2.4) node[above right] {$\eta=0$};
  \fill (6.3,0) circle (1.8pt);
  \node[below] at (6.3,-.05) {$z=0$, \ $v'=-v^3$};
  \node[above] at (3.95,1.9) {leaving};
  \node[above] at (8.65,1.9) {entering};
  \node at (6.3,-.85) {(b)};
\end{tikzpicture}
\caption{The organizing field of \eqref{eq:quadratic-family} for
\(C_0=10\), computed. (a) The Poincar\'e disk, through
\((x,y)\mapsto(x,y)/\sqrt{9+x^2+y^2}\): the graphic \(\Gamma(C_0)\), the
unstable focus at the origin with an orbit spiralling to
\(\Gamma(C_0)\), the section \(\Sigma_0\), and the central section
\(\Sigma_c\) below the nilpotent saddle--node at infinity. (b) The chart
\eqref{eq:projective-coordinates-new}--\eqref{eq:projective-family-new}
near the saddle--node, with the vertical scale enlarged: the line at
infinity \(z=0\), the branches of \(\Gamma(C_0)\) entering the
saddle--node from \(v>0\) and leaving it into \(v<0\), the central
section \(\eta=0\), and orbits in the three sectors of \(z\ge0\).}
\label{fig:organizing-graphic-new}
\end{figure}

\subsection{A fixed neighborhood of the whole graphic}

The Hausdorff condition must force passage through the sections as well
as proximity to the graphic. Removing a neighborhood of the finite
repelling equilibrium gives an isolating block and sections valid for
every nearby parameter.

\begin{proposition}[Isolation and whole-graphic capture]
\label{prop:isolation-capture-new}
There exist a compact set $B_{\mathrm{iso}}\subset\mathbb S^2$, a compact annular
neighborhood $B_{\mathrm{cap}}\Subset\operatorname{int}B_{\mathrm{iso}}$ of $\GammaI$,
a parameter neighborhood $V_{\mathrm{cap}}$, a transverse affine section $\Sigma_0$,
and $\rho_{\mathrm{cap}}>0$ with the following properties.

\begin{enumerate}[label=\textup{(\roman*)}]
\item The maximal invariant subset of $B_{\mathrm{iso}}$ for each field with parameter in
  $V_{\mathrm{cap}}$ is contained in $\operatorname{int}B_{\mathrm{iso}}$.
\item If a periodic orbit $\gamma$ satisfies
  $\distH(\overline\gamma,\GammaI)<\rho_{\mathrm{cap}}$, then
  $\overline\gamma\subset B_{\mathrm{cap}}$ and $\gamma$ crosses
  $\Sigma_0$ and the entrance and exit faces of a fixed projective box
  exactly once per minimal period, in the order of $\GammaI$.
\end{enumerate}
\end{proposition}

A periodic orbit as in (ii) is called \emph{captured}. The reverse
Hausdorff inclusion forces section crossings; in a thin tubular annulus,
their common orientation and the intersection number of an embedded
circle make each crossing unique. Appendix~\ref{app:graphic-capture-details}
gives the neighborhood construction and proof.

\begin{remark}
\label{rem:cyclicity-formulation}
The bound in \cref{thm:direct-finite-cyclicity} also applies after
restricting the orbits to a compact set: for each \(p\in V\) and each
compact \(B\subset\mathbb S^2\), there are at most three limit cycles
\(\gamma\) with \(\overline\gamma\subset B\) and
\(\distH(\overline\gamma,\GammaI)<\varepsilon\).
One such set is \(B_{\mathrm{iso}}\) of
\cref{prop:isolation-capture-new}, whose maximal invariant subset lies
in its interior.
\end{remark}

\subsection{The normal-rate coordinate}

The negative logarithmic rate of \(z\) identifies the central section
and the two directions of passage. Following \cite[Lemma~2.1]{HK2026},
we take this rate, \(\eta=-\dot z/z\), as a coordinate. In the projective time of
\eqref{eq:projective-family-new}, it is
\begin{equation}
\label{eq:normal-rate-coordinate-new}
 \eta=\mu _3+v+v^2+Cvz.
\end{equation}
Since $\partial_v\eta=1+2v+Cz$, this is an analytic coordinate together with
$z$ on a fixed projective box and a sufficiently small parameter
neighborhood.

\begin{lemma}[Normal form]
\label{lem:physical-normal-form-new}
In the coordinates $(\eta,z)$, the projective field has the exact form
\begin{equation}
\label{eq:physical-normal-form-new}
 \dot\eta=q(\eta,p)-zB(\eta,z,p),
 \qquad \dot z=-\eta z,
\end{equation}
where $p=(A,C,\mu _1,\mu _2,\mu _3)$, $B$ and $q$ are analytic, and
\begin{equation}
\label{eq:physical-normal-form-signs-new}
 B>0,\qquad \partial_z\{zB(0,z,p)\}>0
\end{equation}
on a fixed smaller box.  On $z=0$,
\begin{equation}
\label{eq:boundary-drift-conjugacy-new}
 q(\mu _3+v+v^2,p)=(1+2v)P(v).
\end{equation}
At the organizing parameter,
\begin{equation}
\label{eq:organizing-normal-coefficients-new}
 B(0,0)=1,\qquad B_\eta(0,0)=1,\qquad
 q(\eta)=-\eta^3+O(\eta^4).
\end{equation}
Consequently, after one more reduction of the box and parameter
neighborhood, there are positive constants $b_0,b_1,b_2$ such that
\begin{equation}
\label{eq:normal-sign-margins-new}
 B\ge b_0,\qquad B_\eta\ge b_1,\qquad
 3R_q+\eta\partial_\eta R_q\le-b_2,
\end{equation}
where
\begin{equation}
\label{eq:q-expansion-new}
 q(\eta,p)=q_0+q_1\eta+q_2\eta^2+\eta^3R_q(\eta,p).
\end{equation}
\end{lemma}

\begin{proof}
The equation for \(z\) in \eqref{eq:physical-normal-form-new} follows directly
from \eqref{eq:projective-family-new} and
\eqref{eq:normal-rate-coordinate-new}. Differentiating $\eta$ along the
projective vector field \eqref{eq:projective-family-new} gives an analytic first
component.  Its restriction to $z=0$ is
$(1+2v)P(v)$, which proves \eqref{eq:boundary-drift-conjugacy-new}.  The
remainder vanishes on $z=0$ and therefore factors as $-zB$.  Direct
evaluation at the organizing parameter gives
\eqref{eq:organizing-normal-coefficients-new}.  The inequalities follow by
continuity on a compact smaller box.
\end{proof}

For fixed $(A,C)$, the coefficients in \eqref{eq:q-expansion-new} are local
coordinates on $(\mu _1,\mu _2,\mu _3)$, as shown in
Appendix~\ref{app:coordinates}.

\subsection{Regular transitions on fixed flow boxes}

\begin{lemma}[Regular transitions]
\label{lem:regular-transitions}
Fix a finite order \(k\). Let \(\mathcal X_p\) be a \(C^{k+1}\) family
with parameters in a neighborhood of a compact subset of a
finite-dimensional parameter space. Consider transitions between
transverse sections, defined on neighborhoods of the closed section
windows under consideration. Suppose that their orbit segments remain
in a common compact regular region and have flight times bounded by
a fixed \(t_{\max}\). Suppose also that the section parameterizations,
their inverse charts, and their defining functions have uniform
\(C^{k+1}\) bounds, and that the input and output transverse fluxes
have absolute value at least a fixed \(c>0\).
Then the transition maps and their local inverses have uniform \(C^k\)
bounds on the section windows.
\end{lemma}

\begin{proof}
Smooth flow dependence \cite[Theorem~2.11]{Teschl2012}, on finitely
many flow neighborhoods, gives uniform derivatives up to time \(t_{\max}\).
The implicit function theorem for the hitting equation
\(g_p(\varphi_p^t(x))=0\), whose \(t\)-derivative is the nonzero
transverse flux, gives the transition bounds; reversed time gives
the inverse bounds.
\end{proof}


\subsection{The exterior transition}

Fix a small level $z=z_*$.  At the organizing parameter the two branches of
the graphic meet this level in two points with opposite signs of $\eta$.
Let $W_+$ and $W_-$ be compact transverse windows around the positive and
negative endpoints.  The finite part of the graphic defines, for nearby
parameters, a regular first-hit map
\begin{equation}
\label{eq:exterior-map-new}
 \Pi_p:W_-\longrightarrow W_+.
\end{equation}
On both windows we use the auxiliary energy coordinate
\begin{equation}
\label{eq:physical-energy-new}
 I=z_*-\frac{\eta^2}{2}.
\end{equation}
Choose a bounded normal-rate graph box which is downward closed in
the energy variable \(U=\eta^2/2\), for example
\(0\le U\le U_{\max}\), \(0<z\le z_*\), with the organizing
branches away from its upper radius face. It contains the two strips
of \(B_{\mathrm{cap}}\) along the organizing branches with
\(0<z\le z_*\). Choose the organizing endpoints in the interiors
of \(W_\pm\). After reducing the capture data of
\cref{prop:isolation-capture-new} and the parameter neighborhood of
\eqref{eq:exterior-map-new}, every captured periodic orbit stays inside
the graph box, away from its upper radius face, during its central passage
and meets the interiors of both windows. This follows by narrowing
the tubular neighborhood at the two regular section crossings.

\begin{lemma}[Exterior multiplier]
\label{lem:exterior-multiplier-new}
Put
\[
 r_*=\sqrt{2z_*+C_0z_*^2},\qquad a_*=1+C_0z_*,
 \qquad R_*=\frac{r_*}{z_*},
\]
and take $z_*$ small enough that $a_*>2r_*$.  Define the finite-section
transverse multiplier
\[
 m_{\mathrm{aff}}(R_*)=
 \exp\!\left(\int_{-R_*}^{R_*}
       \frac{4x}{x^2+2x+C_0}\,\dd x\right).
\]
Let $I_-^0$ and $I_+^0$ be the two endpoint states of the organizing orbit
in the coordinate \eqref{eq:physical-energy-new}.  The derivative of the
organizing exterior transition at that orbit is
\begin{equation}
\label{eq:exterior-multiplier-new}
 (\Pi_0)'(I_-^0)=
 m_{\mathrm{aff}}(R_*)\frac{(a_*+r_*)(a_*+2r_*)}
                    {(a_*-r_*)(a_*-2r_*)},
 \qquad
 \Lambda_0=\exp\!\left(-\frac{4\pi}{\sqrt{C_0-1}}\right).
\end{equation}
In particular, $(\Pi_0)'(I_-^0)=\Lambda_0\{1+O(\sqrt{z_*})\}$ as
$z_*\downarrow0$.
Let $\Lambda_0<\Lambda_+<1$ and fix $z_*$ sufficiently small. There are
enlarged compact transverse windows \(\widetilde W_\pm\) at \(z=z_*\)
around the organizing endpoints, in whose interiors the windows \(W_\pm\)
are chosen, and \(\Lambda_->0\) such that, for nearby parameters,
\(\Pi_p\) extends to a neighborhood of \(\widetilde W_-\) with values in
the positive section, has uniformly bounded derivatives of every fixed
finite order there, and satisfies
\begin{equation}
\label{eq:exterior-uniform-contraction-new}
 0<\Lambda_-\le\Pi_p'(I)\le\Lambda_+<1\qquad(I\in\widetilde W_-).
\end{equation}
In particular, endpoints uniformly \(o(1)\)-close to \(W_\pm\) lie in
\(\widetilde W_\pm\) for sufficiently small parameters.
\end{lemma}

\begin{proof}
Along the invariant parabola, use $x$ as time and let $\mathfrak k$ be a transverse
variation measured by $H_0$.  Equations
\eqref{eq:parabola-invariance-new}--\eqref{eq:parabola-speed-new} give
\[
 \frac{d}{dx}\log\mathfrak k
 =\frac{4x}{x^2+2x+C_0}.
\]
The finite-section multiplier in this affine transverse coordinate is
$m_{\mathrm{aff}}(R_*)$.  Since
\[
 \PV\int_{-\infty}^{\infty}
 \frac{4x}{x^2+2x+C_0}\,\dd x
 =-\frac{4\pi}{\sqrt{C_0-1}},
\]
and direct integration on the symmetric finite interval gives
\[
 m_{\mathrm{aff}}(R_*)
 =\Lambda_0\{1+O(R_*^{-1})\},
\]
we have $m_{\mathrm{aff}}(R_*)=\Lambda_0\{1+O(\sqrt{z_*})\}$.
We now convert the affine multiplier to the endpoint coordinate \(I\).
At fixed $z=z_*$ the organizing parabola meets the section at $v=\pm r_*$.
Direct differentiation of the
energy coordinate with respect to the affine transverse coordinate $H_0$
gives the two section factors
\[
 z_*^2(a_*+r_*)(a_*+2r_*),\qquad
 z_*^2(a_*-r_*)(a_*-2r_*),
\]
respectively.  Their ratio, multiplied by $m_{\mathrm{aff}}(R_*)$, yields
\eqref{eq:exterior-multiplier-new}; the ratio tends to one as
$z_*\downarrow0$.
Since $\Lambda_0<\Lambda_+$, fix $z_*$ with this derivative below
\(\Lambda_+\). Smooth dependence then gives the asserted extension and
\eqref{eq:exterior-uniform-contraction-new} on sufficiently small
windows \(\widetilde W_\pm\) around the organizing endpoints, for nearby
parameters. The nearby
exterior arcs have uniformly bounded flight times through regular flow
boxes and uniformly transverse endpoints, so
\cref{lem:regular-transitions} gives the finite-order derivative bounds
on \(\widetilde W_-\). The last assertion holds because \(W_\pm\) lie in
the interiors of \(\widetilde W_\pm\).
\end{proof}

\subsection{One central height}

The central section separates the incoming and outgoing halves of the
projective passage. Put
\begin{equation}
\label{eq:central-section-new}
 \Sigma_c(p)=\{\eta=0:0<z<z_*\}.
\end{equation}
The height $h=z|_{\Sigma_c(p)}$ gives a common analytic coordinate on
this parameter-dependent section in the original $(v,z)$ chart.
We call the point
$(\eta,z)=(0,h)$ of $\Sigma_c(p)$ a \emph{central point} and $h$ its
\emph{central height}.

\begin{lemma}[Unique central crossing]
\label{lem:unique-central-crossing-new}
Every periodic orbit captured by \cref{prop:isolation-capture-new} meets
$\Sigma_c(p)$ exactly once during its projective passage, and this
intersection is transverse.
\end{lemma}

\begin{proof}
At $\eta=0$, set
\[
 f_p(z)=q(0,p)-zB(0,z,p).
\]
By \eqref{eq:physical-normal-form-signs-new}, $f_p$ is strictly decreasing.
If $f_p(z)=0$, both components of
\eqref{eq:physical-normal-form-new} vanish, so the point is an equilibrium
and cannot lie on a periodic orbit.  Hence every central crossing of a
captured periodic orbit is transverse.

Suppose a crossing at $z_1$ goes from $\eta<0$ to $\eta>0$.  Until a
hypothetical second crossing, $z$ decreases.  The second crossing would
therefore occur at $z_2<z_1$ and would require $f_p(z_2)<0$, whereas strict
decrease gives $f_p(z_2)>f_p(z_1)>0$.  The opposite orientation is treated
by reversing all inequalities.  There is at most one crossing.  Existence
follows because the entrance and exit faces have opposite signs of $\eta$.
\end{proof}

For the directed passage considered here the crossing is from $\eta>0$ to
$\eta<0$; hence a passing height satisfies
\begin{equation}
\label{eq:central-passing-condition-new}
 f_p(h)<0.
\end{equation}

\subsection{The interval of central returns}

For a passing height $h$, parameterize both half-orbits by increasing
$z$: follow the positive half-orbit backwards and the negative half-orbit
forwards. Put
\begin{equation}
\label{eq:squared-graph-new}
 U_\sigma(z;h)=\frac{\eta_\sigma(z;h)^2}{2},\qquad
 \eta_\sigma=\sigma\sqrt{2U_\sigma},\qquad
 \sigma\in\{+,-\}.
\end{equation}
With $z$ increasing, both functions satisfy
\begin{equation}
\label{eq:physical-graph-equation-new}
 \frac{dU_\sigma}{dz}
 =\mathcal G_{\sigma,p}(U_\sigma,z)
 :=B(\sigma\sqrt{2U_\sigma},z,p)
   -\frac{q(\sigma\sqrt{2U_\sigma},p)}{z},
 \qquad U_\sigma(h;h)=0.
\end{equation}
The smooth field selects a unique branch through the initial
square-root singularity of this coordinate.

\begin{proposition}[Central displacement on one interval]
\label{prop:central-displacement-interval-new}
After reducing the capture data, there is, for every parameter
$p$, an interval $\cH_p\subset(0,z_*)$, possibly empty, such that both graphs
\eqref{eq:physical-graph-equation-new} reach $W_+$ and $W_-$.  The endpoint
functions
\begin{equation}
\label{eq:endpoint-functions-new}
 I_\sigma(h)=z_*-U_\sigma(z_*;h)
\end{equation}
are continuous and strictly increasing on $\cH_p$.  The central first-return
domain is an interval $\cD_p\subset\cH_p$, and its fixed points are exactly
the zeros on $\cH_p$ of
\begin{equation}
\label{eq:central-displacement-new}
 \Delta_p(h)=\Pi_p(I_-(h))-I_+(h).
\end{equation}
For each fixed \(p\), the maps \(I_\pm\) and \(\Delta_p\) are real
analytic on \(\operatorname{int}\cH_p\).
Every captured periodic orbit gives one zero of
\eqref{eq:central-displacement-new}, and every zero gives one periodic
orbit in the selected return neighborhood.  This correspondence preserves
isolation.
\end{proposition}

We call a central height \(h\in\cH_p\) \emph{admissible}, and use the
same adjective for its half-orbits and their endpoints in \(W_\pm\).

\begin{proof}
We first compare the graphs to determine which heights reach the
endpoint windows. We then determine the domain and fixed points of the
first return. Let
\[
 \mathcal Z_p=\{0<h<z_*:f_p(h)<0\}.
\]
This is an interval because $f_p$ is strictly decreasing.  For $h\in
\mathcal Z_p$ and $z>h$, condition
\eqref{eq:central-passing-condition-new} and the strict
decrease of $f_p$ give
\begin{equation}
\label{eq:graph-positive-at-axis-new}
 \mathcal G_{\sigma,p}(0,z)
 =\frac{zB(0,z,p)-q(0,p)}{z}>0.
\end{equation}
Thus a graph which has become positive cannot return to $U=0$ at a larger
height.

Let $h_1<h_2$ be passing heights and suppose first that the graph starting at
$h_1$ exists up to $h_2$.  Then
$U_\sigma(h_2;h_1)>0=U_\sigma(h_2;h_2)$.  Scalar uniqueness away from
$U=0$ gives
\begin{equation}
\label{eq:central-graph-order-new}
 U_\sigma(z;h_1)>U_\sigma(z;h_2)
\end{equation}
for every common later height. If the first graph stays in the fixed
downward-closed graph box until $z_*$, the second stays between it and
zero and also reaches $z_*$ by standard continuation. Consequently, for
each sign, the set $\mathcal Z_{\sigma,p}\subset\mathcal Z_p$ of heights
whose graph reaches $z_*$ inside that box is an interval which
contains, with each of its heights, every larger height of
$\mathcal Z_p$.

To show that \(zB-q>0\) along these graphs, note that initially
\(\mathcal G_{\sigma,p}>0\). At the first zero of
\(\mathcal G_{\sigma,p}(U(z),z)\), if one exists,
one has \(U>0\) and \(q=zB\), so
\begin{equation}
\label{eq:graph-slope-first-zero}
 \frac{d}{dz}\mathcal G_{\sigma,p}(U(z),z)
   =B_z+q/z^2=\frac{B+zB_z}{z}>0.
\end{equation}
This contradicts a first crossing from positive values. The strict
inequality \(B+zB_z>0\) holds on the fixed normal-rate box after
it is reduced.

To determine how the endpoints vary with the central height, differentiate
the moving initial condition at an interior height. This gives
\begin{equation}
\label{eq:moving-height-variation-new}
 \partial_hU_\sigma(z_*;h)
 =\frac{f_p(h)}h\exp\!\left(
  \int_h^{z_*}\partial_U\mathcal G_{\sigma,p}
       (U_\sigma(z;h),z)\,\dd z\right)<0,
\end{equation}
where $f_p(h)/h=q(0,p)/h-B(0,h,p)<0$.
The possible initial $O((z-h)^{-1/2})$ singularity is integrable; equivalently
one may differentiate the smooth flow before passing to $U$.
Hence $U_\sigma(z_*;h)$ is continuous and
strictly decreasing, and $I_\sigma(h)$ is continuous and strictly increasing,
on $\mathcal Z_{\sigma,p}$.

We can now impose the endpoint windows. Each $W_\sigma$ is a compact
interval in the energy coordinate, so
\[
 \mathcal H_{\sigma,p}
 =\{h\in\mathcal Z_{\sigma,p}:I_\sigma(h)\in W_\sigma\}
\]
is an interval, because it is the inverse image of an interval under a
strictly monotone map.  Define
\[
 \mathcal H_p=\mathcal H_{+,p}\cap\mathcal H_{-,p}.
\]
Every captured periodic orbit has its unique central height in
$\mathcal H_p$, by the choice of these strips and of the windows above.
In fact its height belongs to \(\operatorname{int}\mathcal H_p\).
Its central crossing is regular, its endpoint crossings are transverse
and lie in \(\operatorname{int}W_\pm\), and its two compact half-orbits
avoid the radius face. Smooth flow dependence therefore continues
both half-orbits for all central heights in an open neighborhood
of that height, with endpoints still in the two windows.

For fixed \(p\) and \(h_0\in\operatorname{int}\cH_p\), the
analytic field has a regular central point, since \(f_p(h_0)<0\).
The selected half-orbits reach \(z=z_*\) in finite time:
\(\ddot z(0)=-h_0f_p(h_0)>0\) near the central point, and
\(U_\sigma>0\) on the remaining compact graph. At the endpoint,
\(\dot z=-\eta z_*\ne0\). Analytic flow dependence and the analytic
implicit function theorem for these hitting equations therefore give
analytic \(I_\sigma\). The same argument applies to the exterior
transition, whose section coordinate has derivative \(-\eta\ne0\).
Thus \(\Delta_p\) is analytic on \(\operatorname{int}\cH_p\).

It remains to identify the return domain within the admissible interval.
We restrict the central return to heights whose return height lies in
$\cH_p$. The images $I_\sigma(\cH_p)$ are intervals, and this return is defined
precisely when $\Pi_p(I_-(h))$ belongs to $I_+(\cH_p)$. Since $\Pi_p\circ I_-$
is strictly increasing, its inverse image of this interval is an interval
$\cD_p$.  On $\cD_p$ the return is
\[
 \cP_p(h)=I_+^{-1}\bigl(\Pi_p(I_-(h))\bigr).
\]
The equality $\cP_p(h)=h$ is equivalent to
\eqref{eq:central-displacement-new}. Every zero on $\cH_p$ belongs to
$\cD_p$, since its common endpoint lies
in $I_+(\cH_p)$.  Joining the negative half-orbit, the exterior arc, and the
positive half-orbit proves the correspondence with periodic orbits.  The
unique central crossing proves injectivity, and a transverse flow box
preserves isolation in both directions.
\end{proof}

\begin{figure}[t]
\centering
\begin{tikzpicture}[>=Stealth,x=1.05cm,y=.9cm]
  \draw[->,gray!70] (-4.4,0)--(4.4,0) node[right] {$\eta$};
  \draw[->,gray!70] (0,-.2)--(0,4.5) node[above] {$z$};
  \draw[dashed,gray!70] (-3.6,3.8)--(3.6,3.8) node[right] {$z=z_*$};
  \draw[very thick,blue!55!black,
    postaction={decorate},decoration={markings,
      mark=at position .58 with {\arrowreversed{Stealth}}}]
    (0,1) .. controls (1.0,1.25) and (1.8,2.45) .. (2.8,3.8);
  \draw[very thick,blue!55!black,
    postaction={decorate},decoration={markings,
      mark=at position .58 with {\arrow{Stealth}}}]
    (0,1) .. controls (-1.0,1.25) and (-1.8,2.45) .. (-2.8,3.8);
  \draw[very thick,red!65!black,->]
    (-2.8,4.02) to[bend left=28] node[above] {$\Pi_p$} (2.8,4.02);
  \fill (0,1) circle (2pt) node[below right] {$h$};
  \fill (-2.8,3.8) circle (2pt) node[below left] {$I_-(h)$};
  \fill (2.8,3.8) circle (2pt) node[below right] {$I_+(h)$};
  \node[above left] at (-2.8,3.8) {$W_-$};
  \node[above right] at (2.8,3.8) {$W_+$};
  \node[right] at (.08,2.55) {$\Sigma_c(p)$};
\end{tikzpicture}
\caption{The central return in normal-rate coordinates. The
blue curves are the two halves of the orbit through $(0,h)$; the
red arrow is the regular exterior transition.  Positions and curvature are
schematic.  The closing condition is
$\Pi_p(I_-(h))=I_+(h)$.}
\label{fig:central-return-new}
\end{figure}

Returns based at other transverse sections are conjugate to
\(\cP_p\); they have the same number of fixed points, and the
conjugacy preserves isolation.
